\documentclass[11pt]{article}
\usepackage[a4paper,margin=1.1in]{geometry}
\usepackage[T1]{fontenc}
\usepackage{lmodern}
\usepackage{amsmath,amssymb,amsthm}
\usepackage{booktabs}
\usepackage{longtable}
\usepackage{array}
\usepackage{microtype}
\usepackage[hidelinks]{hyperref}
\usepackage{natbib}
\usepackage{setspace}
\usepackage{enumitem}
\usepackage{xcolor}

\newtheorem{theorem}{Theorem}[section]
\newtheorem{proposition}[theorem]{Proposition}
\newtheorem{lemma}[theorem]{Lemma}
\newtheorem{corollary}[theorem]{Corollary}
\theoremstyle{definition}
\newtheorem{definition}[theorem]{Definition}
\newtheorem{assumption}[theorem]{Assumption}
\theoremstyle{remark}
\newtheorem{remark}[theorem]{Remark}

% A formal name, as it appears in the Lean development.
\newcommand{\leanunderscore}{\char`\_\hskip 0pt plus 0.4pt\relax}
\newcommand{\lean}[1]{\begingroup\ttfamily\small\let\_\leanunderscore #1\endgroup}
% The Lean name attached to a stated result.
\newcommand{\formal}[1]{\par\noindent\begingroup\footnotesize\raggedright\emph{Formal statement:} \lean{#1}.\par\endgroup}
\newcommand{\Thorn}{\text{\th}}
\newcommand{\R}{\mathbb{R}}
\newcommand{\E}{\mathbb{E}}
\DeclareMathOperator{\supp}{supp}

\onehalfspacing
\tolerance=600 \emergencystretch=1.5em

\title{Overlapping Generations in Lean:\\ The Life-Cycle Household}
\author{LeanEconomics Org [mostly Claude]\thanks{The Lean development described here is at
\url{https://github.com/LeanEconomics/LeanEconomics}. It was written mostly by Claude (Anthropic),
an AI coding model, working under the direction of Robert Kirkby (Victoria University of
Wellington); every result reported was checked by the Lean kernel. This is the companion to the
write-up on Aiyagari (1994) in the same repository. Comments welcome.}}
\date{September 2026\\[0.5em]\small Draft, in progress}

\begin{document}
\maketitle

\begin{abstract}
\noindent
This is the record of a programme to formalise overlapping-generations economies in Lean 4. A
household that lives $J$ periods against uninsurable earnings risk and a borrowing constraint
solves a finite backward induction, so everything proved for the infinite-horizon Bellman operator
with an arbitrary continuation applies stage by stage with no fixed point. On top of that we prove
the finite-horizon sandwich, $\kappa_k m\le c_k\le\kappa_k(m+H_k(z))$ with
$1/\kappa_k=\sum_{i\le k}(\Thorn/R)^i$, and then that $\kappa_k$ is exactly the reciprocal of that
geometric sum, so the gap to its infinite-horizon limit closes at a known rate. Adding retirement
and a pension, we settle a claim often made informally: a means test does make the policy
non-monotone, but not for the reason usually given. Saving rises with cash on hand whatever the
continuation, by Topkis, so a value function made non-concave by the test cannot by itself reverse
it; what the test reverses is consumption, and cash on hand itself once the taper is steeper than
the gross return. The main result is existence of a stationary equilibrium with both boundary
conditions discharged: capital supply is a finite composition, with no stationary distribution,
no Feller property and no Doeblin condition, and it is bracketed above and below by affine bounds
on a cohort's accumulated assets whose coefficients are explicit recursions in the discount
factor, the interest rate, the earnings process and the horizon. Nothing in the hypotheses refers
to the household's decisions, and at Aiyagari's calibration they hold on an interval of five
percentage points containing the true equilibrium. We also report what the life cycle does for
Light's Theorem 1, that saving rises with the interest rate: it holds at every age below a
critical relative risk aversion of about $3.4$, so the life cycle delivers Aiyagari's $\mu=3$ and
not his $\mu=5$; at logarithmic utility we give a sufficient condition that is true where the
Euler condition is false. Where the bounds stop is the same obstruction at both ends, that a power
mean with negative exponent is superadditive rather than $1$-Lipschitz.
\end{abstract}

\noindent\emph{JEL:} C63, D52, E21. \quad \emph{Keywords:} overlapping generations, life cycle,
Bewley models, formal verification, Lean.

\section{Introduction}

The companion write-up in this repository proved that the stationary equilibrium of
\citet{Aiyagari1994} is unique at his logarithmic calibration, and located exactly what stands in
the way for relative risk aversion $\mu\in\{3,5\}$. Every obstacle it met traces to the fixed point
of the infinite-horizon problem: a condition on the response of consumption to the interest rate
fails in a tail of wealth that the stationary distribution reaches; an induction that would prove
the response bound has a step multiplier $1/\Thorn>1$, so it cannot close as a supremum argument;
and the only floor on capital supply comes from an ergodic identity that loses a factor
$\|u'\|^2/\operatorname{Var}$. Overlapping generations remove all three at once. A household
that lives $J$ periods solves a finite backward induction from an explicit terminal condition, so
a multiplier $1/\Thorn$ per step compounds to $\Thorn^{-J}$, which at Aiyagari's numbers is about
$1.03$. Wealth is bounded by $J$ periods of accumulation, so any condition needs checking on a
known bounded region. And the distribution of the youngest cohort is common to every interest
rate, so an ordering of policies pushes forward to an ordering of the whole age profile of
distributions by induction on age, with no Doeblin atom and no uniqueness of a stationary
distribution.

The target of the programme is uniqueness of the stationary equilibrium of a stochastic life-cycle
economy in the manner of \citet{Huggett1996} at $\mu\in\{3,5\}$. It is reached at $\mu=1$ and not
above, and the paper is as much a record of where it is blocked as of what is proved.

Sections~\ref{sec:diamond} to~\ref{sec:eqm} are the foundations. Section~\ref{sec:diamond} does
Diamond's two-period economy, where the steady state is unique for every risk aversion and the
household-side comparative static that blocks the general case runs the wrong way without doing any
harm. Section~\ref{sec:lc} sets out the life-cycle household, the finite-horizon sandwich, and the
propensity floor. Section~\ref{sec:ret} adds retirement, a pension and a means test, and settles
which policy a means test makes non-monotone and in which state space. Section~\ref{sec:eqm} proves
existence of a stationary equilibrium with both boundary conditions discharged from primitives rather
than assumed: capital supply in a life cycle is a finite composition of finite sums, so existence
needs no stationary distribution and no compactness argument beyond Berge, and it can be bracketed by
two affine bounds on a cohort's accumulated assets sharing a slope, with every coefficient an explicit
recursion in the primitives. At Aiyagari's calibration the hypotheses hold on an interval of about
five percentage points containing the true rate of $5.09\%$.

Sections~\ref{sec:thm1} to~\ref{sec:inc} are about uniqueness, and they are best read as three
attempts on one obstacle. Uniqueness needs capital supply to rise with the interest rate, and the
only route the literature offers is Light's Theorem 1, that saving rises with the rate at every
state. Section~\ref{sec:thm1} reports how far that goes: it holds at every age below a critical
relative risk aversion of about $3.4$, so the life cycle covers Aiyagari's $\mu=3$ and not his
$\mu=5$; the Euler reduction of it is false at logarithmic utility exactly where an equilibrium sits,
and a different sufficient condition, on the value gap between two rates, is true there instead
(Section~\ref{sec:log}). The section also separates the mechanism from the usual story: the reversal
happens in the certainty-equivalent household, which faces no risk at all, so precautionary saving is
not what drives it, and we give the criterion in closed form and prove its certainty-equivalent half
for every risk aversion.

Section~\ref{sec:mfg} asks what mean field games offer. The Lasry--Lions argument was dismissed for
the stationary Bewley model because it needs a common initial distribution, which overlapping
generations have; so the monotonicity condition becomes a question to check rather than one to
dismiss. Date by date it fails, at every risk aversion, and not for the expected reason. What the
argument actually uses is weaker --- only the summed pairing is signed --- and that weaker condition
holds throughout the bracket in which the equilibrium is known to lie. But pushed to its weakest form
it collapses to monotone capital supply, so the route is not logically weaker than the one it was
meant to bypass. The section is a closure, with the reduction and the two theorems it produces kept
because they are about the life-cycle household rather than about mean field games.

Section~\ref{sec:inc} is the route that works at $\mu=1$. Single crossing does not ask capital supply
to rise, only to fall more slowly than the firm's demand; over the proved bracket the demand curve is
steep enough that a proof may lose more than the true increment and still close. Coupling the two
economies on a common earnings path and splitting the step at the lower-rate state gives an envelope
in which the failure of Theorem 1 is priced rather than excluded, and its per-stage multiplier is
damped by the propensity floor. Two further facts make the estimate small enough: a cohort's earnings
marginal is the invariant law at every age, so the mass weighting is exact in the coordinate where the
failure lives; and the multiplier is one plus the variance of the risk tilt, an identity rather than a
bound. That variance is certified below what the damping allows, in exact rational arithmetic, at every
stage and over the whole asset range. What that delivers is Theorem 1 at every state and age, hence
monotone capital supply, hence uniqueness --- at logarithmic utility, at one calibration, with the last
hypothesis discharged by certificate rather than by proof term. Section~\ref{sec:prog} records what
remains.

As in the companion, every result stated with a formal name has been checked by the Lean kernel
against Mathlib; the development is at
\url{https://github.com/LeanEconomics/LeanEconomics}, directory \texttt{LeanEconomics/OLG}, and
the numerical work is in \texttt{WriteUps/WriteUpOLG/numerics}.

\section{The life-cycle household}\label{sec:lc}

\subsection{Stages}

The household of the companion write-up is kept: earnings $y(z)$ follow a finite Markov chain
$\pi$, the gross return is $R=1+r$, cash on hand is $m=Ra+y(z)$, saving $a'\in[0,\bar a]$ is
bounded below by the borrowing limit and above by an artefactual cap, utility is CRRA with
$u'(c)=c^{-\mu}$ and $\beta\in(0,1)$. What changes is the horizon: the household lives $J$
periods and leaves nothing behind. Count periods from the end. With $k$ periods still to come
after today, the continuation value is the $k$-th Bellman iterate from zero,
\[
  V_k=T^k 0,\qquad (Tv)(a,z)=\max_{a'}\ \bigl\{u(m-a')+\beta\textstyle\sum_{z'}\pi(z,z')v(a',z')\bigr\},
\]
and the household with $k$ periods to come saves $g_k=\operatorname{policyOf}V_k$ and consumes
$c_k=m-g_k$ (\lean{stageValue}, \lean{stagePolicy}, \lean{stageConsumption}). Age $j$ of a
$J$-period life is stage $k=J-1-j$. In the last period the household saves nothing and consumes
its cash on hand (\lean{stagePolicy\_zero}, \lean{stageConsumption\_zero}).

This is the finite-horizon problem with a flat age profile of earnings: the same income process at
every age. An age profile, retirement and survival risk need income to be a parameter of the
Bellman step rather than a field of the structure, and are deferred to the next stage of the
programme.

\subsection{What transfers without a fixed point}

The companion development proved its household results for the Bellman operator applied to an
arbitrary continuation $v$, and only then passed to the fixed point. Here there is no fixed point
to pass to, and each result is read at $v=V_k$:

\begin{itemize}[leftmargin=1.5em]
\item consumption is positive at every stage when marginal utility is unbounded at zero
  (\lean{stageConsumption\_pos});
\item the stage value has concave slices, so the optimal action is unique and both saving and
  consumption rise with wealth (\lean{stagePolicy\_mono}, \lean{stageConsumption\_mono});
\item Carroll--Kimball: consumption is concave in wealth at every stage
  (\lean{concaveOn\_stageConsumption});
\item the Euler inequality under the cap, $\beta R\,\E[u'(c_k(a',z'))\mid z]\le u'(c_{k+1}(a,z))$,
  at every state (\lean{stage\_euler\_le}), and the Euler inequality above the floor, the reverse,
  at every state where the household saves (\lean{stage\_euler\_ge}).
\end{itemize}

\subsection{The finite-horizon sandwich}

Let $\Thorn=(\beta R)^{1/\mu}$ and define the finite-horizon minimal propensities to consume by
\[
  \kappa_0=1,\qquad \kappa_{k+1}=\frac{\kappa_kR}{\Thorn+\kappa_kR},\qquad\text{so that}\qquad
  \frac1{\kappa_k}=\sum_{i=0}^{k}\Bigl(\frac{\Thorn}{R}\Bigr)^{i}
\]
(\lean{stageMPC}): the share of cash on hand consumed by a household with $k$ periods to come and
no future income. Let $H_0=0$ and $RH_{k+1}(z)=\sum_{z'}\pi(z,z')\bigl(y(z')+H_k(z')\bigr)$, the
expected present value of the next $k$ incomes (\lean{stageHumanWealth}).

\begin{theorem}[The finite-horizon sandwich]\label{thm:sandwich}
Assume CRRA utility with marginal utility unbounded at zero, and that the cap is slack at every
stage. Then for every $k$, every $z$ and every $a\in[0,\bar a]$,
\[
  \kappa_k\,m\ \le\ c_k(a,z)\ \le\ \kappa_k\bigl(m+H_k(z)\bigr).
\]
\formal{IncomeFluctuation.stageMPC\_mul\_le\_stageConsumption,
IncomeFluctuation.stageConsumption\_le\_stageHumanWealth}
\end{theorem}

Each bound is one Euler step, applied $k$ times. The lower bound: if consumption against $v$ is at
least $\kappa m$, the Euler inequality above the floor at a state where the household saves gives
$c^{-\mu}\le\beta R\,\E[(\kappa Ra')^{-\mu}]$ with $a'=m-c$, hence $\kappa R(m-c)\le\Thorn c$ and
$c\ge\kappa R(\Thorn+\kappa R)^{-1}m$; at a corner $c=m$ (\lean{crra\_stageMPC\_step}). The upper
bound: if consumption against $v$ is at most $\kappa(m+H(z))$, the Euler inequality under the cap
and Jensen for the negative power give $c^{-\mu}\ge\beta R\,(\E[c'])^{-\mu}$ with
$\E[c']\le\kappa R(a'+H'(z))$, hence $c\le\kappa R(\Thorn+\kappa R)^{-1}(m+H'(z))$
(\lean{crra\_stageUpper\_step}). The recursion for $\kappa$ is the same on both sides, which is
why the sandwich is exact in the homogeneous case.

The propensities decrease in the horizon towards the infinite-horizon $\kappa=1-\Thorn/R$ of the
companion write-up: $\kappa\le\kappa_k$ for every $k$ (\lean{minMPC\_le\_stageMPC}), so a
household with less life left consumes a larger share of its cash, and the infinite-horizon bounds
are the limiting case. A first consequence for the programme is an explicit bound on saving,
$g_k(a,z)\le(1-\kappa_k)(Ra+y_{\max})$ (\lean{stagePolicy\_le}), which iterated from zero wealth
at the start of life bounds the wealth of every cohort.

\section{Retirement, a pension, and the means test}\label{sec:ret}

\subsection{Retirement with a lump-sum pension}

The household works until age $J_r$ and is retired after it, receiving a pension in place of its
labour earnings. With a \emph{lump-sum} pension $p$ the retired household faces the same problem
as the working one with income replaced by the constant $p$, so the retired phase is itself an
income-fluctuation problem (\lean{withPension}) and everything in Section~\ref{sec:lc} applies to
it unchanged. The life-cycle value is assembled by iterating the working Bellman operator on the
retired value: $\mathrm{lifeValue}(n_r,n_w)$ is the value of a household with $n_w$ working and
$n_r$ retired periods still to come (\lean{lifeValue}), so a household of age $j<J_r$ in a life of
$J$ periods sits at $n_w=J_r-j$ and $n_r=J-J_r$. Saving and consumption rise with wealth at every
age (\lean{lifePolicy\_mono}, \lean{lifeConsumption\_mono}).

Two facts about the pension itself are worth recording. A retiree's human wealth is the same in
every earnings state, because the pension does not depend on the earnings shock
(\lean{stageHumanWealth\_withPension\_const}), and it never exceeds the perpetuity value of the
pension, $H_k\le p/r$ (\lean{stageHumanWealth\_withPension\_le}). With Theorem~\ref{thm:sandwich}
the retiree's consumption therefore satisfies
\[
  \kappa_k\,m\ \le\ c_k(a)\ \le\ \kappa_k\Bigl(m+\frac{p}{r}\Bigr)
\]
at every retired age. And cash on hand is strictly increasing in assets
(\lean{withPension\_resources\_strictMono}). Both of these fail once the pension is means-tested.

\subsection{Means testing}

A pension is means-tested when the amount paid falls with the recipient's assets. We write the
schedule as a function $p(\cdot)$ of assets and cash on hand as $m(a)=Ra+p(a)$ (\lean{cash}),
and consider three schedules (\lean{lumpSum}, \lean{tapered}, \lean{cliff}): a lump sum; a taper,
withdrawn at rate $\tau$ per unit of assets above a threshold and never negative; and a cliff,
withdrawn entirely above the threshold.

It is often argued that a means test makes the household's policy function non-monotone, and the
reason usually given is that the test makes the value function non-concave. The argument as
usually given is wrong, and the conclusion is nevertheless true, for a different reason. Three
results separate the two.

\begin{theorem}[Monotonicity in cash on hand is never lost]\label{thm:topkis}
Let period utility be strictly concave on the positive half-line and let $W$ be \emph{any}
continuation value. If $m_1<m_2$ and $x_i$ is optimal at $m_i$, then $x_1\le x_2$.
\formal{OLG.isOptimalSaving\_mono}
\end{theorem}

The proof is Topkis. Adding the two optimality inequalities gives
$u(m_1-x_1)+u(m_2-x_2)\ge u(m_1-x_2)+u(m_2-x_1)$; if $x_2<x_1$ the two right-hand consumptions are
strict convex combinations of the two left-hand ones with the same sum, so strict concavity makes
the reverse inequality strict, a contradiction. The continuation enters only through terms that
cancel. So a non-concave value function --- which a means test certainly produces --- cannot by
itself make saving fall with cash on hand, whatever the shape of $W$.

What a non-concave continuation does produce is a jump, and the jump falls on consumption.

\begin{theorem}[Consumption falls when wealth rises]\label{thm:cons}
There are a continuation value arising from a means-tested pension and two levels of cash on hand,
the second larger, at which optimal consumption is strictly smaller. With log utility, a unit
gross return, and a pension of one withdrawn entirely above assets of one: at $m=6$ the household
saves exactly $1$ and consumes $5$; at $m=8$ it saves $4$ and consumes $4$.
\formal{OLG.exists\_consumption\_decrease}
\end{theorem}

At $m=6$ the household keeps the pension, saving the most it can without losing it; at $m=8$ the
pension is no longer worth keeping, so the household gives it up and saves the proceeds. Saving
rises from $1$ to $4$, as Theorem~\ref{thm:topkis} says it must, and since it rises by more than
wealth does, consumption falls. Both optima are exact: the objective is $\log$ of a quadratic on
each side of the threshold, and the comparison is checked by the kernel.

The third result is the genuine non-monotonicity of the saving policy, and it comes from the
budget rather than from the value function. A retiree's pension is assessed each period on the
assets then held, so the schedule enters cash on hand directly.

\begin{theorem}[A gentle taper is harmless; a cliff is not]\label{thm:taper}
If $0\le\tau<R$ then $a\mapsto Ra+p(a)$ is strictly increasing under the taper
(\lean{cash\_strictMono\_tapered}), and so, by Theorem~\ref{thm:topkis}, saving is monotone in
assets (\lean{isOptimalSaving\_mono\_assets}). Under a cliff it is not: if $a\le\bar a<b$ and
$R(b-a)<p$ then cash on hand is strictly smaller at $b$ (\lean{cash\_lt\_cash\_of\_cliff}), and
there are two asset levels, the second larger, at which the household saves strictly less: with a
concave continuation, at assets $1$ it saves $1/2$ and at assets $3/2$ only $1/4$.
\formal{OLG.exists\_saving\_decrease\_in\_assets}
\end{theorem}

The continuation in the second half of Theorem~\ref{thm:taper} is $\log(x+1)$, which is concave:
nothing is hidden in it. The non-monotonicity is entirely the budget's, and the threshold is
sharp, $\tau$ against $R$. A taper gentler than the gross return leaves the household's saving
monotone in assets no matter how badly the test distorts the value function; a taper steeper than
the gross return, or a cliff, makes a richer household poorer in the only sense that matters for
today's decision.

\subsection{Which policy, and in which variable}

The two non-monotonicity results are about different objects, and it is worth separating them in
the household's own state space, assets $a$ and the Markov income shock $z$.

Next period's assets are monotone. For each $z$, the optimal $a'$ is nondecreasing in $a$
whatever the continuation value, with no concavity assumed of it:

\begin{theorem}[$a'$ rises with $a$, at each $z$]\label{thm:azmono}
Let $a<a'$ both lie in the asset region. Then for every continuation $v$ and every $z$,
$g_v(a,z)\le g_v(a',z)$.
\formal{IncomeFluctuation.policyOf\_mono\_of\_lt}
\end{theorem}

This is Theorem~\ref{thm:topkis} read in the state space: cash on hand rises with assets, so
optimal saving does. The repository already had this statement (\lean{policyOf\_mono}) but with
concave slices of the continuation as a hypothesis, which is exactly what a means test destroys;
the Topkis argument removes the hypothesis, at the cost of requiring the two asset levels to be
distinct. The continuation cancels between the two optimality inequalities and only strict
concavity of period utility survives.

Consumption is not monotone, and the contrast is sharp rather than an artefact of proof
technique. Consumption is $m(a)-g_v(a,z)$, and both terms rise with $a$; when the continuation is
non-concave the second can rise faster, which is exactly Theorem~\ref{thm:cons}. So monotone
consumption genuinely requires concavity of the continuation, while monotone saving does not.

Two things follow for the means test. Where the pension enters only next period's problem, so
that today's cash on hand is unaffected, Theorem~\ref{thm:azmono} applies verbatim: next period's
assets still rise with this period's, and only consumption misbehaves. Where the test is assessed
on the assets currently held, the hypothesis of Theorem~\ref{thm:azmono} is the one that fails:
cash on hand is no longer increasing in $a$, and with it the policy. That is the whole of
Theorem~\ref{thm:taper}, and the dividing line is $\tau$ against $R$.

\subsection{At realistic parameters}

The three theorems predict exactly which of the two channels operates for which schedule, and a
calibrated retiree confirms it. We solve a deterministic retiree of $25$ years by backward
induction with global search over a grid of $6001$ asset levels --- the value function is not
concave, so a first-order condition cannot be used --- at $r=4\%$, $\beta=0.96$, relative risk
aversion $3$, a full pension of $0.40$ of mean earnings and a free area of $3$ mean-earnings units
(\texttt{numerics/meanstest.m}). The taper of $0.078$ per unit of assets is Australia's asset test,
three dollars a fortnight per thousand dollars.

\begingroup\small
\begin{center}
\begin{tabular}{@{}lccc@{}}
\toprule
Pension schedule & Cash on hand & Saving & Consumption \\
\midrule
Lump sum & increasing & increasing & increasing \\
Taper $0.078$ (Australia) & increasing & increasing & falls by $0.025$--$0.046$ \\
Taper $0.50$ & increasing & increasing & falls by $0.13$--$0.24$ \\
Taper $1.20$ (above $R$) & falls & falls by $0.004$ & falls by $0.55$--$0.70$ \\
Cliff & falls by $0.40$ & falls by $0.32$--$0.37$ & falls by $0.55$--$0.70$ \\
\bottomrule
\end{tabular}
\end{center}
\endgroup

Every entry is what the theorems require. Saving falls with assets in exactly the two rows where
cash on hand does, which are exactly the two rows whose taper exceeds the gross return, as
Theorem~\ref{thm:taper} says; in the other three it rises, as Theorem~\ref{thm:topkis} says it
must. Consumption falls in every row with a test of any kind, as Theorem~\ref{thm:cons} says it
can, and the drop sits exactly at the free area and grows with age, from $0.025$ of mean earnings
at the start of retirement to $0.046$ twenty years in under the Australian taper. The mechanism is
the one the witness isolates: just above the free area the effective return to saving is only
$R-\tau$ until the pension is exhausted, so the household jumps its saving toward the far side of
the taper range and cuts consumption to pay for it.

The effect is not confined to retirees. A worker in the last year before retirement, facing the
Australian taper for the whole of retirement, has a monotone saving policy and a consumption
policy that falls by $0.025$ of mean earnings at the same threshold: the non-concavity is
inherited through the continuation value, and it is again consumption that carries it.

\section{Stationary equilibrium}\label{sec:eqm}

Every cohort is born with no assets and an earnings state drawn from a fixed law $\nu$, lives $J$
periods, and is replaced; cohorts have equal weight. The distribution of the youngest cohort is
therefore the same at every interest rate, and the distribution of every older cohort is that law
pushed forward by the policies of the ages in between. This is the structural advantage of the
life cycle over the infinite horizon, and it removes most of the machinery the stationary case
needed.

Two of the consumption bounds this section uses are developed in Section~\ref{sec:thm1}, where
they arise in the hunt for the critical risk aversion: the constraint-incidence bound of
Theorem~\ref{thm:incidence} and the asset-dependent saving floor of Section~\ref{sec:assetdep}.
Nothing here depends on the verdict of that section.

\subsection{The aggregate is a finite composition}

Because assets at birth are degenerate at zero, no measure theory is required. Let
\[
  (\mathcal{S}_k h)(a,z)=\sum_{z'}\pi(z,z')\,h\bigl(g_k(a,z),z'\bigr)
\]
be the expectation of $h$ next period for a household with $k$ periods still to come
(\lean{stageStep}). The assets a household holds over the rest of its life, summed across ages
and seen from its current state, satisfy
\[
  A_0(a,z)=a,\qquad A_{k+1}(a,z)=a+(\mathcal{S}_{k+1}A_k)(a,z)
\]
(\lean{cohortAssets}), and aggregate capital per head is
$K=\frac1J\sum_z\nu(z)\,A_{J-1}(0,z)$ (\lean{olgCapital}). There is no pushforward measure, no
Feller property, no stationary distribution, no Doeblin condition and no uniqueness of an
invariant law: the equilibrium object is a finite sum of finite sums, and the recursion above is
its definition.

\subsection{The chain}

\begin{theorem}[Light's Theorem 2 for the life cycle]\label{thm:thm2}
Suppose two economies share the earnings chain and their policies are ordered at every age,
$g^{(1)}_k\le g^{(2)}_k$ on the asset region. Then $A^{(1)}_k\le A^{(2)}_k$ for every $k$, and
aggregate capital is ordered, $K^{(1)}\le K^{(2)}$.
\formal{IncomeFluctuation.cohortAssets\_le,
IncomeFluctuation.olgCapital\_le\_of\_stagePolicy\_le}
\end{theorem}

The proof is an induction on age with two ingredients, and both are already available. The class
of functions increasing in assets is preserved by $\mathcal{S}_k$
(\lean{monoRegion\_stageStep}), because next period's assets rise with this period's at each
earnings state and for any continuation, which is Theorem~\ref{thm:azmono}. And ordered policies
keep ordered functions ordered (\lean{stageStep\_le\_stageStep}): the higher-rate economy
evaluates a larger function at a larger asset level. Note what is \emph{not} needed. The
infinite-horizon version of this step required monotone transitions, a coupling argument and
first-order stochastic dominance of the stationary distributions; here the whole statement is an
inequality between two functions on the same finite state space, because the distribution the two
economies start from is literally the same.

\subsection{Uniqueness}

\begin{theorem}[Single crossing]\label{thm:thm3}
A capital supply that never falls with the rate meets a strictly falling demand at most once.
\formal{IncomeFluctuation.olgEquilibriumRate\_unique}
\end{theorem}

With Theorem~\ref{thm:thm2} this reduces uniqueness of the stationary equilibrium to Light's
Theorem 1, that saving rises with the rate at every age, which is the subject of
Section~\ref{sec:thm1}: it holds at Aiyagari's calibration for relative risk aversion up to about
$3.4$, so the chain closes at his $\mu=3$ and not at his $\mu=5$.

\subsection{Continuity, and existence}

The last input is continuity of capital supply in the rate. In the infinite horizon this was a
substantial development, because the stationary distribution itself had to be shown to move
continuously. Here there is no stationary distribution: capital supply is a finite composition of
finitely many policies, so all that is needed is that each policy moves continuously with the
rate, and that continuity composes.

We take the route of the companion write-up and put the interest rate into the state, so that
Berge's theorem gives continuity in assets and the rate jointly rather than continuity in a
parameter. That device was used there at the fixed point, but everything it rests on --- that the
feasible sets and the objectives agree on the slice --- holds for an arbitrary continuation. So
it works stage by stage, and here it is easier, because the slice of the augmented Bellman
iterate is the iterate of the sliced programme by a plain induction, with no fixed point to
identify (\lean{sliceAt\_bellman}, \lean{sliceAt\_augStageValue}). The augmented argmax at each
stage is then a singleton, since the sliced one is by Carroll--Kimball concavity of the iterates
(\lean{argmax\_aug\_stage\_eq\_singleton}), so the stage policy is jointly continuous
(\lean{continuousOn\_augStagePolicy}) and agrees with the sliced one
(\lean{augStagePolicy\_eq}). Running the cohort recursion in the augmented state
(\lean{augCohortAssets}) gives a composition of continuous maps that reads, at a rate in range, as
capital supply at that rate (\lean{augCohortAssets\_eq}, \lean{olgSupply\_eq}).

\begin{theorem}[Capital supply is continuous]\label{thm:cont}
$r\mapsto K(r)$ is continuous on $[r_{lo},r_{hi}]$.
\formal{IncomeFluctuation.continuousOn\_olgSupply}
\end{theorem}

\begin{theorem}[Existence]\label{thm:exist}
If capital supply lies below demand at the bottom of the interval and above it at the top, and
demand is continuous, then the two meet: there is a rate in $[r_{lo},r_{hi}]$ at which the capital
the cohorts accumulate equals the capital the firm demands.
\formal{IncomeFluctuation.exists\_olgEquilibrium}
\end{theorem}

Both hypotheses of Theorem~\ref{thm:exist} are about the solved model. The rest of this section
replaces them with inequalities that never mention the household's policy or value function.

\subsection{Where the household can be}\label{sec:reach}

Proposition~\ref{prop:floor} inherits from Theorem~\ref{thm:sandwich} the hypothesis that the
artefactual cap is slack, and the convenient way to state that is slack everywhere in
$[0,\bar a]$. That hypothesis is not innocent, and it fails exactly where the floor is wanted.

The one-step bound gives $g_k(a,z)\le(1-\kappa_k)(Ra+y_{\max})$, so $[0,\bar a]$ is carried into
itself only if $(1-\kappa_k)(R\bar a+y_{\max})\le\bar a$, which needs $(1-\kappa_k)R<1$. As
$k$ grows $\kappa_k$ falls to $1-\text{\th}/R$, so the requirement is $\text{\th}<1$, that is
$\beta R<1$: the household must be impatient. Above the rate of time preference no cap is forward
invariant at any level, because a patient household sitting at the cap wants to save past it. The
cap then binds, the Euler inequality under it is unavailable, and the slackness hypothesis is
false rather than merely unverified. A theorem carrying it says nothing there.

This matters because a stationary life-cycle equilibrium sits above the rate of time preference.
In the calibration of Section~\ref{sec:num2} with $\lambda=1/\beta-1=4.17\%$, capital supply
crosses demand near $5\%$, and the rate at which the \emph{provable} floor crosses demand is
near $9\%$. Both are in the region where no cap is slack everywhere.

A cohort that begins life with nothing never visits the top of $[0,\bar a]$. The development
therefore carries a family of regions, one per stage, that the stage policies map down through,
and asks for slackness only on those.

\begin{definition}[Reachable family]\label{def:regions}
A reachable family assigns to each stage $k$ a set of asset levels such that every set sits
inside $[0,\bar a]$, the stage-$(k+1)$ policy carries the stage-$(k+1)$ set into the stage-$k$
set, and $g_k<\bar a$ on every set.
\formal{IncomeFluctuation.StageRegions}
\end{definition}

The whole interval $[0,\bar a]$ is a reachable family exactly when the cap is slack on it
(\lean{allRegions}), so nothing is lost: Theorem~\ref{thm:sandwich} and
Proposition~\ref{prop:floor} are the constant-family instances of statements that hold on any
reachable family (\lean{stageConsumption\_le\_stageHumanWealth\_on}, \lean{le\_stagePolicy\_on}).

What a cohort visits is bounded by feasibility alone. Set $\rho_0=0$ and
$\rho_{j+1}=R\rho_j+y_{\max}$: a household that cannot carry forward more than its cash on hand
holds at most $\rho_j$ after $j$ periods of life, whatever the earnings draws.

\begin{proposition}[A cohort's own regions]\label{prop:reach}
If $\rho_{K+1}<\bar a$ then the sets $[0,\rho_{K-k}]$, empty for $k>K$, form a reachable family
for a cohort of length $K+1$ that starts life with nothing.
\formal{IncomeFluctuation.reachRegions}
\end{proposition}

The bound is deliberately crude, using none of the damping that optimal saving supplies, and that
is what keeps it free of the sandwich it is used to prove. Its cost is a larger cap. At $r=10\%$
over sixty periods $\rho_{K+1}$ is about $9{,}044$ where the optimal path reaches about $82$.
Since the cap is an artefact of the state space rather than an economic object, a cap of ten
thousand is as acceptable as a cap of one hundred, and unlike slackness everywhere it is a
condition that can actually be met. The sandwich and the floor then hold along a cohort's path
with no hypothesis on the cap beyond $\rho_{K+1}<\bar a$
(\lean{stageConsumption\_le\_stageHumanWealth\_reach}, \lean{le\_stagePolicy\_reach}).

At the bottom of the rate interval the same crude bound does the other half of the work on its
own. When $R<1$ the sequence $\rho_j$ converges to $y_{\max}/(1-R)$, so cohort assets are bounded
without any appeal to optimality: at $r=-7\%$ the bound is $42.1$ against a capital demand of
$56.3$.

\subsection{The propensity in closed form}\label{sec:geom}

Both bounds below are built from the finite-horizon propensity $\kappa_k$ and from human wealth,
and both need to know not just where $\kappa_k$ sits but how fast it settles.

\begin{proposition}[The finite-horizon propensity]\label{prop:kappa}
Write $q=\Thorn/R$. Then $\kappa_k\sum_{i\le k}q^i=1$, and hence
$\kappa_k\bigl(1-q^{k+1}\bigr)=1-q$: the propensity exceeds its infinite-horizon limit $1-q$ by
exactly $\kappa_kq^{k+1}$.
\formal{IncomeFluctuation.stageMPC\_mul\_mpcSum},
\lean{IncomeFluctuation.stageMPC\_mul\_one\_sub\_pow}
\end{proposition}

The defining recursion $\kappa_{k+1}=\kappa_kR/(\Thorn+\kappa_kR)$ is exactly
$1/\kappa_{k+1}=1+q/\kappa_k$, so the reciprocal accumulates the geometric series and nothing
else. The same sum tames the slope that both bounds share. Writing $S_k=\sum_{i\le k}q^i$, the
recursion $G_0=1$, $G_{k+1}=1+G_k(1-\kappa_{k+1})R$ has an age-dependent coefficient, but
$W_k=G_kS_k$ obeys $W_{k+1}=S_{k+1}+\Thorn W_k$, which does not. Then $S_{k+1}\ge1$ and
$(1-q)S_k\le1$ give the bound that matters.

\begin{proposition}[A slope that grows]\label{prop:slope}
$G_k\ \ge\ (1-q)\sum_{j\le k}\Thorn^{\,j}$.
\formal{IncomeFluctuation.mul\_geomSum\_le\_cohortSlope}
\end{proposition}

Where the crude $G_k\ge1$ does not grow at all, this grows geometrically, and
Section~\ref{sec:sharpnum} reports what that is worth.

\subsection{A floor on capital}\label{sec:clampfloor}

Existence needs capital supply above demand at the top of the rate interval. In the infinite
horizon the only floor available came from an ergodic identity that loses a factor
$\|u'\|_\infty^2/\operatorname{Var}$. The life cycle needs nothing of the kind: the upper sandwich
of Theorem~\ref{thm:sandwich} bounds consumption at every age, and what is not consumed is saved.

\begin{proposition}[Affine floor on saving]\label{prop:floor}
At every age and every state in the region,
\[
  g_k(a,z)\ \ge\ (1-\kappa_k)\bigl(y_z+Ra\bigr)-\kappa_k H_k(z).
\]
\formal{IncomeFluctuation.le\_stagePolicy\_on}, \lean{IncomeFluctuation.le\_stagePolicy}
\end{proposition}

Iterating it is not a matter of tracking a distribution over assets. An affine lower bound on a
cohort's accumulated assets suffices, $F_k(z)+G_ka\le A_k(a,z)$, with the slope $G$ of
Section~\ref{sec:geom} common to all earnings states and an intercept carrying one value per
state:
\[
  F_0=0,\qquad
  F_{k+1}(z)=\textstyle\sum_{z'}\pi(z,z')F_k(z')
    +G_k\bigl[(1-\kappa_{k+1})y_z-\kappa_{k+1}H_{k+1}(z)\bigr].
\]

\begin{proposition}[Affine floor on accumulated assets]\label{prop:affine}
On any reachable family, $F_k(z)+G_ka\le A_k(a,z)$; hence capital per head is at least
$\sum_z\nu_zF_K(z)/(K+1)$.
\formal{IncomeFluctuation.cohortFloor\_add\_mul\_le\_cohortAssets},
\lean{IncomeFluctuation.le\_olgCapital}
\end{proposition}

The state dependence of $F$ is the whole content, not a refinement. Reading
Proposition~\ref{prop:floor} at the worst earnings state instead, which is the only way to get a
bound uniform in $z$, gives a floor that is identically zero at every rate in the calibration of
Section~\ref{sec:num2}, against $4.75$ at $r=10\%$ for the state-dependent one.

Averaging against weights the chain leaves alone collapses the vector recursion to a scalar one,
and the sign of each increment then has an exact answer.

\begin{theorem}[The floor rises exactly when the household is patient]\label{thm:patient}
Let $\bar y$ and $\bar H_k$ be the averages of earnings and of human wealth under weights $\nu$
satisfying $\nu\pi=\nu$. If $\Thorn\ge1$, equivalently $\beta R\ge1$, then
$\kappa_k\bar H_k\le(1-\kappa_k)\bar y$ at every age, so the aggregate floor never falls and is
never negative.
\formal{IncomeFluctuation.stageMPC\_mul\_sum\_stageHumanWealth\_le},
\lean{IncomeFluctuation.sum\_mul\_cohortFloor\_mono}
\end{theorem}

The increment of the aggregate floor is $G_k\kappa_{k+1}\bar y$ times the slack in that
inequality, and the inequality reverses below $\Thorn=1$. So the provable floor is flat under the
rate of time preference and climbs above it, which is why the upper endpoint of the rate interval
has to be taken above $\lambda=1/\beta-1$. The proof is one induction: writing
$v_k=(1-\kappa_k)/\kappa_k$ and $u_k=\bar H_k/\bar y$, both satisfy first-order recursions,
$v_{k+1}=(\Thorn/R)(1+v_k)$ and $u_{k+1}=(1+u_k)/R$, from the same starting point, and
$\Thorn\ge1$ keeps the first above the second.

One thing the affine form cannot do is refuse to claim negative assets. Its intercept goes badly
negative: at the calibration, $160$ of $413$ entries are below zero and the worst is $-62$, with
$F_K$ running from $-58.8$ at the lowest earnings to $+685$ at the highest. A cohort's accumulated
assets cannot be negative, so all of that is slack the aggregate inherits. Reading the floor as a
function of assets rather than as an affine form fixes it, because the saving floor can then be
clamped at zero where it belongs:
\[
  \Lambda_0(z,a)=a,\qquad
  \Lambda_{k+1}(z,a)=a+\textstyle\sum_{z'}\pi(z,z')\,\Lambda_k\bigl(z',\Phi_{k+1}(z,a)\bigr),
\]
with $\Phi$ the clamped saving floor of Section~\ref{sec:assetdep}. The proof needs only that
$\Lambda_k$ is increasing in assets, which the clamp preserves.

\begin{proposition}[The clamped floor]\label{prop:clampfloor}
$\Lambda_k(z,a)\le A_k(a,z)$ on any reachable family, so capital per head is at least
$\sum_z\nu_z\Lambda_K(z,0)/(K+1)$, which is at least what Proposition~\ref{prop:affine} gives.
\formal{IncomeFluctuation.cohortFloorAt\_le\_cohortAssets},
\lean{IncomeFluctuation.le\_olgCapital\_at}
\end{proposition}

\subsection{A ceiling on capital}\label{sec:affceil}

Existence also needs capital supply below demand at the bottom of the interval, and there three
bounds are available, each asking for more and giving more.

The cheapest uses the budget and nothing else. A household cannot carry forward more than its cash
on hand, so an asset level $B$ with $RB+y_{\max}\le B$ is never left once entered.

\begin{proposition}[A forward-invariant level caps supply]\label{prop:ceiling}
If $0\le B\le\bar a$ and $RB+y_{\max}\le B$, then capital per head is at most $B$.
\formal{IncomeFluctuation.olgCapital\_le\_of\_invariant}
\end{proposition}

The second says the household will not want to carry that much. Consuming at least
$\kappa_k(m+H_k(z))$ by the constraint-incidence bound of Theorem~\ref{thm:incidence}, it saves at
most $(1-\kappa_k)(y_z+Ra)-\kappa_kH_k(z)$.

\begin{proposition}[The incidence ceiling]\label{prop:ceiling2}
If $B\le\bar a$ and, at every age and earnings state,
$(1-\kappa_j)(y_w+RB)-\kappa_jH_j(w)\le B$, then capital per head is at most $B$.
\formal{IncomeFluctuation.olgCapital\_le\_of\_incidence}
\end{proposition}

Both bound \emph{every} age's assets by one level, and so pay $(K+1)B$ for a cohort that holds
nothing when it is born, nothing when it dies, and anything like $B$ only in between. The floor
never made that mistake, and the third ceiling does not either: it is the same affine object,
built from the same slope $G$, with the saving ceiling in place of the saving floor,
\[
  \bar F_0=0,\qquad
  \bar F_{k+1}(z)=\textstyle\sum_{z'}\pi(z,z')\bar F_k(z')
    +G_k\bigl[(1-\kappa_{k+1})y_z-\kappa_{k+1}H_{k+1}(z)\bigr],
\]
with $H$ here the capped, risk-adjusted human wealth of Theorem~\ref{thm:incidence}.

\begin{proposition}[The affine ceiling]\label{prop:affceil}
$A_k(a,z)\le\bar F_k(z)+G_ka$, and hence capital per head is at most
$\sum_z\nu_z\bar F_K(z)/(K+1)$.
\formal{IncomeFluctuation.cohortAssets\_le\_cohortCeil},
\lean{IncomeFluctuation.olgCapital\_le\_cohortCeil}
\end{proposition}

Replacing a worst case over ages by an average over them is worth a factor of fourteen, and it
keeps working above the rate of time preference, where no forward-invariant level exists at all.

\begingroup\small
\begin{center}
\begin{tabular}{@{}lccccc@{}}
\toprule
& $r=-7\%$ & $r=-4\%$ & $r=0\%$ & $r=\lambda$ & $r=5\%$\\
\midrule
Capital demand & $56.25$ & $14.06$ & $7.03$ & $4.62$ & $4.33$\\
Proposition~\ref{prop:ceiling}, the budget & $42.52$ & $74.41$ & --- & --- & ---\\
Proposition~\ref{prop:ceiling2}, incidence & $22.17$ & $35.11$ & --- & --- & ---\\
Proposition~\ref{prop:affceil}, affine & $1.54$ & $2.28$ & $3.86$ & $5.98$ & $6.47$\\
True capital supply & $0.64$ & $1.03$ & $1.98$ & $3.78$ & $4.25$\\
\bottomrule
\end{tabular}
\end{center}
\endgroup

\subsection{Existence from primitives}

The two ends can now be put together, and the statement has taken four forms as the bounds
improved. All four are in the development; the last is the one to read.

\begin{theorem}[Existence from primitives]\label{thm:exist2}
Fix a rate interval and continuous demand $D$. Suppose there is a $B$ with $0\le B\le\bar a$,
$(1+r_{lo})B+y_{\max}\le B$ and $B\le D(r_{lo})$; and suppose $\rho_{K+1}<\bar a$ and
$D(r_{hi})\le\sum_z\nu_zF_K(z)/(K+1)$ at $r_{hi}$. Then there is a rate in the interval at which
the capital the cohorts accumulate equals the capital the firm demands.
\formal{IncomeFluctuation.exists\_olgEquilibrium\_rate},
\lean{IncomeFluctuation.exists\_olgEquilibrium\_of\_bounds}
\end{theorem}

Nothing in those hypotheses refers to the household's decisions. They are inequalities among the
discount factor, the interest rate, the earnings process, the horizon and the demand curve, with
$\rho$, $F$ and $G$ explicit finite recursions in those primitives. Three refinements follow.
Proposition~\ref{prop:slope} turns the floor into a closed form with no recursion left to run
(\lean{IncomeFluctuation.exists\_olgEquilibrium\_of\_geomBounds}, Theorem~\ref{thm:exist3} below);
and Proposition~\ref{prop:ceiling2} replaces the budget ceiling
(\lean{IncomeFluctuation.exists\_olgEquilibrium\_of\_sharpBounds}). The affine ceiling replaces it
again, and at that point no auxiliary $B$ and no invariance condition survive.

\begin{theorem}[Existence, both ends read the same way]\label{thm:exist4}
With the ceiling at $r_{lo}$ and the floor at $r_{hi}$ both taken as weighted averages of affine
bounds on a cohort's accumulated assets, a stationary equilibrium exists.
\formal{IncomeFluctuation.exists\_olgEquilibrium\_affine}
\end{theorem}

Reading the floor at the assets held as well gives the tightest form.

\begin{theorem}[Existence with both ends read at the assets held]\label{thm:exist5}
With the affine ceiling at $r_{lo}$ and the clamped floor at $r_{hi}$, a stationary equilibrium
exists.
\formal{IncomeFluctuation.exists\_olgEquilibrium\_clamped}
\end{theorem}

The two boundary conditions are then the same object built twice, once from the saving ceiling and
once from the saving floor, sharing a slope. Each half asks for what it can have: at the bottom
the artefactual cap is slack, which holds below the rate of time preference; at the top it cannot
be, and the reachable family of Section~\ref{sec:reach} stands in. So the low endpoint must sit
below $\lambda$ and the high endpoint above it, and Section~\ref{sec:sharpnum} shows a bracket that
does.

For completeness, the closed-form variant.

\begin{theorem}[Existence in closed form]\label{thm:exist3}
Under the conditions of Theorem~\ref{thm:exist2}, with the floor hypothesis replaced by the
closed-form expression of Proposition~\ref{prop:slope} and with $1\le\Thorn\le R$ at $r_{hi}$, a
stationary equilibrium exists.
\formal{IncomeFluctuation.exists\_olgEquilibrium\_of\_geomBounds}
\end{theorem}

\subsection{What the numbers say}\label{sec:num2}

The calibration is the one used throughout: $J=60$, $\beta=0.96$, relative risk aversion $3$, the
Tauchen discretisation of Aiyagari's earnings process with $\rho=0.9$ and $\sigma=0.4$, and a
Cobb--Douglas technology with $\alpha=9/25$ and $\delta=2/25$, so capital demand is
$\alpha/((1-\alpha)(r+\delta))$ and $\lambda=4.17\%$. Capital supply is the average asset holding
across the sixty cohorts, each born with nothing. The last row is not the solved supply but the
floor of Proposition~\ref{prop:affine} (\texttt{numerics/olgexist.m}).

\begingroup\small
\begin{center}
\begin{tabular}{@{}lccccccc@{}}
\toprule
$r$ & $-4\%$ & $0\%$ & $\lambda$ & $6\%$ & $8\%$ & $10\%$ & $15\%$ \\
\midrule
Capital demand & $14.06$ & $7.03$ & $4.62$ & $4.02$ & $3.52$ & $3.13$ & $2.45$ \\
Capital supply & $1.03$ & $1.98$ & $3.78$ & $4.85$ & $6.18$ & $7.69$ & $12.49$ \\
Provable floor & $0.00$ & $0.00$ & $0.00$ & $1.52$ & $3.13$ & $4.75$ & $9.40$ \\
\bottomrule
\end{tabular}
\end{center}
\endgroup

Supply crosses demand at $5.09\%$, above the rate of time preference, which is what
Section~\ref{sec:reach} has to accommodate; and the floor is flat below $\lambda$ and climbs above
it, which is Theorem~\ref{thm:patient}. The optimal path reaches assets of $82$ at $r=10\%$,
against the crude reachable bound of $9{,}033$; the constants are ordinary, unlike the stationary
case, where the corresponding instantiation put $r_{hi}$ within $10^{-27}$ of $\lambda$.

\subsection{The theorem checked, hypothesis by hypothesis}\label{sec:sharpnum}

Take $r_{lo}=2\%$, $r_{hi}=8\%$ and an artefactual cap of $10^4$. Every hypothesis of
Theorem~\ref{thm:exist5} holds (\texttt{numerics/clampcheck.m}): the cap is forward invariant at
$r_{lo}$ with margin $134$; the ceiling is $4.81$ against a demand of $5.63$; a lifetime of
feasible saving reaches $3{,}730$ against the cap; and the floor is $3.80$ against a demand of
$3.52$. The bracket contains the true equilibrium at $5.09\%$.

Sweeping the endpoints shows what each bound is worth. The low endpoint may be taken as high as

\begingroup\small
\begin{center}
\begin{tabular}{@{}lc@{\qquad}lc@{}}
\toprule
Ceiling used & highest $r_{lo}$ & Floor used & lowest $r_{hi}$\\
\midrule
Proposition~\ref{prop:ceiling}, the budget & $-6.73\%$
  & closed form, crude $G_k\ge1$ & $610\%$\\
Proposition~\ref{prop:ceiling2}, incidence & $-5.85\%$
  & closed form, Proposition~\ref{prop:slope} & $18.2\%$\\
Proposition~\ref{prop:affceil}, affine & $+2.78\%$
  & Proposition~\ref{prop:affine}, affine & $8.38\%$\\
& & Proposition~\ref{prop:clampfloor}, clamped & $7.67\%$\\
\bottomrule
\end{tabular}
\end{center}
\endgroup

So the tightest bracket the development gives is $[2.78\%,\,7.67\%]$, and the absolute limit is
$5.09\%$ from both sides, since no bound on capital supply can cross demand before supply does.
The closed-form floor is reported because it needs no recursion at all; at a calibrated rate it is
much the weaker, giving $0.48$ at $r=10\%$ against a true floor of $4.75$.

Neither bracket is tight, and the same fact stops them at both ends. Appendix~\ref{app:routes} gives
it: the obstruction at the top and the obstruction at the bottom are one obstruction, that a power
mean with negative exponent is superadditive rather than $1$-Lipschitz.

\section{Theorem 1 in the life cycle: how far it goes}\label{sec:thm1}

The programme of Section~\ref{sec:prog} rests on Light's Theorem 1, that saving rises with the
interest rate at every state. In the infinite horizon this is open above logarithmic utility. The
life cycle was the reason for hope: a finite backward induction from an explicit terminal
condition, with the per-step loss compounding only $J$ times instead of to a fixed point. This
section reports what the numbers say, at $\beta=0.96$, a seven-state Tauchen chain with
$\rho=0.9$ and $\sigma=0.4$, a zero borrowing limit, $J=60$, and a flat earnings profile
(\texttt{numerics/lc\_check.m}, \texttt{lc\_robust.m}, \texttt{lc\_induct.m},
\texttt{lc\_verify.m}, \texttt{lc\_bound.m}, \texttt{lc\_crit.m}).

\subsection{The statement is its own terminal condition}

The budget constraint makes two statements identical. Since $a'=Ra+y(z)-c$ at either rate,
\[
  g_k(a,z;r_2)\ \ge\ g_k(a,z;r_1)
  \qquad\Longleftrightarrow\qquad
  c_k(a,z;r_2)-c_k(a,z;r_1)\ \le\ (r_2-r_1)\,a ,
\]
so Theorem 1 \emph{is} the affine bound on the rate response of consumption with slope one and
zero intercept. That bound holds with equality in the last period of life, where $c_0=y(z)+Ra$.
The question is therefore whether the terminal condition propagates backward, and it is a
question about one inequality rather than about a pair of constants to be chosen.

\subsection{It propagates, up to a critical risk aversion}

At $\mu=1$ and $\mu=3$ it propagates exactly. Over all sixty stages and the whole wealth region
reachable from zero wealth at birth, the largest violation of $g_k^{(1)}\le g_k^{(2)}$ is zero,
and the consumption bound holds to five parts in $10^{14}$; refining the asset grid from $2000$
to $16000$ points changes nothing. At $\mu=5$ it fails. The failure is a genuine feature of the
household's problem, not of the solution method: the violation is identical at every grid size,
and endogenous-grid and value-function iteration with a pure grid search agree on it to the
resolution of the grid. Saving at zero wealth in the highest income state, in the first period of
a sixty-period life, falls as the rate rises:

\begingroup\small
\begin{center}
\begin{tabular}{@{}lccccc@{}}
\toprule
$r$ & $2\%$ & $3\%$ & $4\%$ & $5\%$ & $6\%$ \\
\midrule
$g(0,z_{\max})$, endogenous grid & $1.7783$ & $1.7638$ & $1.7452$ & $1.7234$ & $1.6993$ \\
$g(0,z_{\max})$, value iteration & $1.7806$ & $1.7606$ & $1.7406$ & $1.7206$ & $1.7006$ \\
\bottomrule
\end{tabular}
\end{center}
\endgroup

Between the two lies a critical degree of risk aversion, and it is sharp. Bisecting on $\mu$ for
the largest violation over all stages gives:

\begingroup\small
\begin{center}
\begin{tabular}{@{}lccc@{}}
\toprule
Horizon $J$ & $r=2\%$ & $r=4\%$ & $r=\lambda=4.17\%$ \\
\midrule
$20$ & $3.79$ & $3.53$ & $3.51$ \\
$40$ & $3.77$ & $3.39$ & $3.36$ \\
$60$ & $3.77$ & $3.39$ & $3.36$ \\
\bottomrule
\end{tabular}
\end{center}
\endgroup

Theorem 1 holds throughout the life cycle below the critical value and fails above it. The value
is stable in the horizon beyond about forty periods, and falls as the rate rises towards the rate
of time preference. Aiyagari's $\mu=3$ is inside it with little room, at $\mu=3.39$ against
$3$; his $\mu=5$ is well outside. So the life cycle buys one of the two open cases and not the
other, and it buys $\mu=3$ only at rates below the rate of time preference, which is where the
equilibrium lies.

\subsection{Why the natural induction does not close}

The backward induction that the transfer lemma suggests compares the two Euler equations at a
state where the lower-rate household saves. Writing $b=g^{(1)}_{k+1}(a,z)$ for the pivot, the
induction hypothesis and monotonicity of $c_k$ give
$c^{(2)}_k(b')\le c^{(1)}_k(b)+\Delta r\,b$ at the relevant points, and the step closes provided
\[
  R_2\,\E\bigl[(c^{(1)}_k(b)+\Delta r\,b)^{-\mu}\bigr]\ \ge\
  R_1\,\E\bigl[c^{(1)}_k(b)^{-\mu}\bigr]\Bigl(1+\frac{\Delta r\,a}{c^{(1)}_{k+1}(a,z)}\Bigr)^{-\mu}.
\]
To first order this asks $\mu\bigl(b/c^{(1)}_k(b)-a/c^{(1)}_{k+1}(a,z)\bigr)\le1/R$: the ratio of
assets to consumption may rise from one age to the next by at most $1/(\mu R)$. It does not. At
zero wealth in the highest income state the left-hand side is large, because such a household
saves a great deal against a certain fall in income, while the right-hand side is about one. The
margin is negative from the third stage of life onward for every degree of risk aversion we tried,
logarithmic utility included, at $-0.12$ for $\mu=1$, $-2.86$ for $\mu=3$ and $-6.29$ for $\mu=5$
per unit of the rate gap. Keeping the term the step discards, the marginal propensity to consume
times the size of the hypothetical violation, bounds the violation by a multiple of $\Delta r$
rather than eliminating it, which is not enough: the aggregate difference the equilibrium argument
must sign is itself of order $\Delta r$.

\subsection{The reduction, formalised, and its verdict}

The backward induction above is proved in Lean, in the form the life cycle makes natural.
Because the budget constraint makes the two statements identical
(\lean{stagePolicy\_le\_iff\_consumption\_le}), the terminal stage needs no hypothesis at all: in
the last period both households consume their cash on hand and save nothing
(\lean{stagePolicy\_mono\_zero}). Corners need none either: if the lower-rate household saves
nothing there is nothing to violate. What the step needs is one inequality per age and state,
which we call the \emph{age-indexed slack condition}: writing $b$ for the lower-rate household's
saving at that state and $C'(z')$ for its consumption next period,
\begin{equation}\label{eq:stageslack}
  R_1\sum_{z'}\pi(z,z')C'(z')^{-\mu}\ \le\
  R_2\sum_{z'}\pi(z,z')\bigl(C'(z')+\Delta r\,b\bigr)^{-\mu}
  \Bigl(1+\frac{\Delta r\,a}{c}\Bigr)^{\mu}
\end{equation}
(\lean{StageSlack}).

\begin{theorem}[The reduction]\label{thm:reduction}
If \eqref{eq:stageslack} holds at every age and every state in the region, then saving rises with
the rate at every age, and aggregate capital rises with the rate.
\formal{IncomeFluctuation.stagePolicy\_mono\_step,
IncomeFluctuation.stagePolicy\_mono\_of\_stageSlack,
IncomeFluctuation.olgCapital\_le\_of\_stageSlack}
\end{theorem}

With Theorem~\ref{thm:thm3} this closes the chain: \eqref{eq:stageslack} at every age gives a
unique equilibrium rate. The question is therefore whether \eqref{eq:stageslack} is true, and the
answer is the one reported above, now stated sharply. Condition \eqref{eq:stageslack} is exactly
the margin tabulated in the previous subsection. At $\mu=3$ it holds for the last three or four
periods of life, with margins $+0.68$, $+0.42$ and $+0.18$ per unit of the rate gap, and fails
from about the fifth-to-last, at $-0.23$, deteriorating to $-2.86$ at birth. At $\mu=1$ it holds
for most of the life and still fails for the youngest ages, at $-0.12$ at birth.

That last observation is the important one. The condition fails even at logarithmic utility,
where Theorem 1 is certainly true. So the failure is not a symptom of high risk aversion: the
pointwise Euler comparison is simply the wrong instrument in a life cycle, and no sharpening of
it that we could find repairs the young ages. Keeping the term the step discards, the marginal
propensity to consume times the size of the hypothetical violation, only bounds the violation by
a multiple of $\Delta r$, and both refinements we tried --- an upper bound on the propensity to
consume from concavity and the sandwich, and a self-improving bootstrap on the size of the
violation --- converge to a positive bound rather than to zero.

So $\mu=3$ is not delivered by this route, and $\mu=1$ is not either, although the companion
write-up proves the latter by the rescaling argument, which is an induction on Bellman iterates
and so carries over to the life cycle unchanged. What the reduction does deliver is the end of
life, and the exact inequality that any proof must beat.

So the state of the programme is this. Theorem 1 is true in the life cycle for $\mu=3$, and the
statement to prove has the cleanest possible form, its own terminal condition propagated; the
reduction to one inequality per age is proved; and that inequality is false for the young, at
every degree of risk aversion. What is missing is an argument at zero wealth that does not go
through a pointwise comparison of the two Euler equations. The same zero-wealth state is where the
infinite-horizon slack condition was tightest, which suggests the two problems share a single
missing ingredient rather than two.

\subsection{It is not precautionary saving}

A household at zero wealth in the highest earnings state saves heavily, and the obvious diagnosis
is precautionary: it is insuring against the fall in earnings it expects. If that were the
mechanism, precautionary-saving theory would be the place to look for the missing argument. It is
not the mechanism. The same reversal happens in the \emph{certainty-equivalent} household, which
faces the conditional mean earnings path and no risk at all, and there everything is in closed
form (\texttt{numerics/precaut2.m}, \texttt{durcheck.m}).

\begingroup\small
\begin{center}
\begin{tabular}{@{}lcccc@{}}
\toprule
$\mathrm{d}(\text{saving})/\mathrm{d}r$ at zero wealth & $\mu=1$ & $\mu=3$ & $\mu=5$ & $\mu=8$\\
\midrule
Certainty equivalent (mean path, no risk) & $+21.3$ & $+2.81$ & $-0.85$ & $-2.91$\\
With risk & $+17.1$ & $+0.92$ & $-2.05$ & $-3.43$\\
\bottomrule
\end{tabular}
\end{center}
\endgroup

The sign flips between $\mu=3$ and $\mu=5$ in both columns. Bisecting, the certainty-equivalent
threshold is $\mu=4.33$ against $\mu=3.39$ with risk: risk moves the threshold by about one unit
of risk aversion and does not create the effect.

Without risk the mechanism is exact. A household with $J$ periods, no borrowing constraint and
CRRA utility has $c_{t+1}=(\beta R)^{1/\mu}c_t$ and a lifetime budget, so
\[
  c_0=\frac{W(R)}{D(R)},\qquad
  W(R)=\sum_t y_tR^{-t},\qquad
  D(R)=\sum_t\beta^{t/\mu}R^{t/\mu-t}
\]
(\lean{consumption\_mul\_pathPrice}), and saving at zero wealth rises with the rate exactly when
$W(R_2)D(R_1)\le W(R_1)D(R_2)$ (\lean{consumption\_antitone\_iff}). Both terms fall with the rate,
$W$ because future earnings are discounted harder and $D$ because the optimal path gets cheaper,
and the comparison is between the two elasticities. Differentiating,
\[
  \frac{\mathrm{d}\log c_0}{\mathrm{d}\log R}=-T_y+\Bigl(1-\frac1\mu\Bigr)T_c ,
\]
with $T_y$ and $T_c$ the Macaulay durations of the two streams, so $c_0$ falls, and saving rises,
exactly when $T_y>(1-1/\mu)T_c$. The predicted elasticity matches the exact one to three decimals
at every $\mu$ we checked. At Aiyagari's numbers $T_y=14.31$ and $T_c\approx18.6$, nearly
independent of $\mu$, so the condition is $1-1/\mu<0.769$, that is $\mu<4.33$.

The economics is then plain, and has nothing to do with risk. A higher interest rate discounts
future earnings more heavily, which makes the household poorer and pushes consumption down; it
also makes a rising consumption path cheaper, which pushes consumption today down only to the
extent that the household is willing to tilt its path. A more risk averse household will not
tilt, so the second force weakens as $\mu$ grows while the first does not, and past $\mu\approx4$
the household responds to a higher rate by consuming more today. The income stream simply has to
be longer than the consumption weights, discounted by the strength of the tilt.

\begin{theorem}[No risk, and risk aversion at most one]\label{thm:ce}
If $\mu\le1$ the comparison holds term by term, whatever the earnings profile: the
certainty-equivalent household always saves more at a higher rate.
\formal{CertaintyEquivalent.humanWealth\_mul\_pathPrice\_le\_of\_le\_one,
CertaintyEquivalent.consumption\_antitone\_of\_le\_one}
\end{theorem}

For $\mu>1$ it is the duration comparison, which is a finite computation in the primitives. So a
precautionary argument near zero wealth would be aimed at the wrong target: what has to be
controlled is the elasticity of human wealth against the elasticity of the price of the
consumption path, and risk enters only as a correction to it.

\subsection{The criterion in elasticity units, and the size of the risk correction}

The certainty-equivalent criterion is cleanest when written in the intertemporal elasticity of
substitution rather than in risk aversion. Since it is linear in $1/\mu$,
\[
  T_y>\Bigl(1-\frac1\mu\Bigr)T_c
  \qquad\Longleftrightarrow\qquad
  \frac1\mu\ >\ 1-\frac{T_y}{T_c} ,
\]
so the elasticity of substitution must exceed one minus the ratio of the two durations. At
Aiyagari's numbers that right-hand side is $0.2306$, giving $\mu<4.33$ as before, and the identity
holds to four decimals against the bisected threshold.

With risk the threshold moves down, and the move is what has to be bounded. Write
\[
  \Delta\ :=\ \frac1{\mu^{\text{risky}}}-\frac1{\mu^{\text{CE}}}
\]
for the gap in elasticity units and $s=\sigma\sqrt{1-\rho^2}$ for the conditional standard
deviation of log earnings, the dispersion the household actually faces one period ahead. Measuring
both thresholds at the same state over fifteen calibrations, $\rho\in\{0,0.6,0.9\}$ and
$\sigma$ from $0.025$ to $0.4$ (\texttt{numerics/gapscale.m}, \texttt{gapfit.m}):

\begingroup\small
\begin{center}
\begin{tabular}{@{}lcccccc@{}}
\toprule
$\rho=0.9$, $\sigma$ & $0.025$ & $0.05$ & $0.1$ & $0.2$ & $0.4$ & fitted exponent\\
\midrule
$\Delta$ & $0.0034$ & $0.0070$ & $0.0144$ & $0.0300$ & $0.0612$ & \\
$\Delta/s$ & $0.315$ & $0.320$ & $0.330$ & $0.344$ & $0.351$ & $1.04$\\
\bottomrule
\end{tabular}
\end{center}
\endgroup

The gap is first order in the conditional standard deviation, not second order: fitting
$\Delta\sim s^{b}$ gives $b=1.04$ at $\rho=0.9$, $0.89$ at $\rho=0.6$ and $0.93$ at $\rho=0$, and
the ratio $\Delta/s$ stays between $0.196$ and $0.351$ across every calibration we ran. That
first-order behaviour is itself informative: a precautionary effect on saving would be second
order in the dispersion, whereas the probability that the borrowing constraint binds in a
low-earnings state next period moves at first order, and that is our reading of where the
correction comes from.

Taking the largest ratio observed as the bound, $\Delta\le0.36\,s$, gives a criterion entirely in
the primitives:
\begin{equation}\label{eq:practical}
  \frac1\mu\ \ge\ 1-\frac{T_y}{T_c}\ +\ 0.36\,\sigma\sqrt{1-\rho^2}.
\end{equation}
At Aiyagari's calibration the right-hand side is $0.2306+0.0628=0.2934$, that is $\mu\le3.41$,
against the $3.39$ that bisection on the full life-cycle model gives. The bound therefore
reproduces the measured threshold to within two per cent, and it says the same thing the earlier
sections said: the calibration he used for $\mu=3$ is inside, and $\mu=5$ is not.

\subsubsection*{The certainty-equivalent half, for every risk aversion}

The first inequality of \eqref{eq:practical} was a theorem only for $\mu\le1$, where the exponent
$\theta=1-1/\mu$ is nonpositive and every factor moves the same way. It is now a theorem for every
$\mu$. Writing $\lambda=R_1/R_2<1$, the criterion $W(R_2)D(R_1)\le W(R_1)D(R_2)$ compares how the
same weights respond to $\lambda^t$ against $(\lambda^t)^\theta$. Each $\lambda^t$ lies in
$[\lambda^{J-1},1]$, and on that interval the concave $x\mapsto x^\theta$ lies above its chord,
which is affine; replacing the power by its chord leaves a comparison of two sums linear in
$\lambda^t$.

\begin{theorem}[The chord criterion]\label{thm:chord}
For $\mu>1$, if
$W(R_2)D(R_1)\le W(R_1)\bigl(\iota\,D(R_1)+\varsigma\,D_\lambda(R_1)\bigr)$, where $\iota$ and
$\varsigma$ are the intercept and slope of that chord and $D_\lambda$ is the path price with each
term discounted once more by $\lambda^t$, then initial consumption falls with the rate.
\formal{CertaintyEquivalent.consumption\_antitone\_of\_chord},
\lean{CertaintyEquivalent.chord\_le\_pathPrice}
\end{theorem}

The replacement costs almost nothing, because it is first-order exact: as $R_2\to R_1$ it
reproduces $T_y\ge(1-1/\mu)T_c$ exactly. At Aiyagari's numbers the chord threshold is $4.3263$
against the exact $4.3267$, and the cost grows only with the size of the rate step, to $0.04$ at
$\Delta r=10^{-3}$ and $0.40$ at $10^{-2}$ (\texttt{numerics/chord.m}). Together with the $\mu\le1$
theorem this settles the certainty-equivalent household at every risk aversion.

\subsubsection*{The risk half, and why it resists}

Two caveats remain. The constant $0.36$ is measured over fifteen calibrations and
is not proved. And both are statements at zero wealth in the highest earnings state, which is
where the life-cycle model binds; the reduction of Theorem~\ref{thm:reduction} is what would carry
them to the rest of the state space, and its own condition fails, so this is a sharp description
of where the truth lies rather than a proof of it.

It is worth being precise about why the constant resists. What has to be bounded is a difference of
\emph{rate responses} between two problems, the risky one and its certainty equivalent, and the
threshold is the rate at which one of those responses changes sign. No bound on the \emph{level} of
consumption helps with that, however sharp: at Aiyagari's calibration the sharpest sandwich
available has width $0.30$ to $0.48$ at the state in question, while the response itself runs from
$-2.48$ at $r=2\%$ through zero to $+0.29$ at $r=6\%$ (\texttt{numerics/scale.m}). A bound on where
consumption lies says nothing about how it moves. What is needed is a bound on the rate response
itself, and that is exactly the object the companion write-up proved cannot be had from the affine
induction, whose per-step multiplier is at least $1/\Thorn>1$. The two halves of
\eqref{eq:practical} therefore stand on different ground: the first is now a theorem for every
risk aversion, and the second is blocked by the same obstruction that closed the rate-response
bound.

\subsubsection*{How much slack a bound on the correction has}

That said, ``blocked'' understates how good the position is, and the quantity worth measuring is
not the correction but the room around it. Write the risky household's saving response at the state
as the certainty-equivalent response less a correction,
$\partial g/\partial r=\partial g_{\mathrm{CE}}/\partial r-\mathrm{corr}$. Theorem 1 at the state
follows from any bound $B\ge\mathrm{corr}$ with $B\le\partial g_{\mathrm{CE}}/\partial r$
(\lean{CertaintyEquivalent.saving\_nonneg\_of\_correction\_le}), so what matters is the ratio
$(\partial g_{\mathrm{CE}}/\partial r)/\mathrm{corr}$: a bound loose by less than that factor still
delivers the theorem. Measured at $r=4\%$, zero wealth, top earnings state:

\begin{center}
\begin{tabular}{lrrrrrrr}
\toprule
$\gamma$ & $1$ & $2$ & $2.5$ & $3$ & $3.2$ & $3.39$ & $5$\\
\midrule
$\partial g_{\mathrm{CE}}/\partial r$ & $20.98$ & $7.28$ & $4.55$ & $2.73$ & $2.17$ & $1.69$ & $-0.89$\\
$\partial g/\partial r$ (risky) & $17.08$ & $4.88$ & $2.49$ & $0.92$ & $0.44$ & $0.03$ & $-2.05$\\
correction & $3.89$ & $2.40$ & $2.06$ & $1.82$ & $1.73$ & $1.66$ & $1.15$\\
\midrule
slack factor & $\mathbf{5.4}$ & $\mathbf{3.0}$ & $\mathbf{2.2}$ & $\mathbf{1.5}$ & $1.25$ & $1.02$ & ---\\
\bottomrule
\end{tabular}
\end{center}

Two things stand out. The correction is \emph{stable}: it stays between $1.15$ and $3.89$ across
$\gamma\in[1,5]$, and it \emph{falls} as risk aversion rises, while the certainty-equivalent response
collapses from $20.98$ to below zero. The threshold is therefore set by the certainty-equivalent side
and not by risk at all --- which is the same conclusion the closed form gave, now with the correction
measured rather than inferred. And the slack is generous where it matters: a bound on the correction
may be wrong by a factor of $5.4$ at $\gamma=1$, $3.0$ at $\gamma=2$ and $1.5$ at $\gamma=3$ and
still certify Theorem 1 there.

That is a far weaker demand than the one left open in Section~\ref{sec:inc}, where the remaining
estimate needs an absolute error below $0.03$ on a quantity of size about one. Here a crude
precautionary bound would do. Note only that the slack is \emph{not} uniform in $\gamma$: the
correction peaks at $3.89$ at $\gamma=1$ while $\partial g_{\mathrm{CE}}/\partial r$ is only $2.73$
at $\gamma=3$, so a single constant cannot cover the range --- the bound must be allowed to depend on
$\gamma$, which costs nothing, since Theorem 1 is wanted at a calibration and not as a family.
Numerics: \texttt{riskgap.m}.

\subsubsection*{The precautionary split, and what a bound on it must achieve}

The slack above is measured on derivatives at $r=4\%$. Theorem 1 is a statement about an interval, so
it is worth redoing in secant form, where the object to be bounded becomes concrete. Write the risky
household's saving at the state as the certainty-equivalent saving plus precautionary saving,
$g=g_{\mathrm{CE}}+P$ with $P=c_0^{\mathrm{CE}}-c_0\ge0$. Then Theorem 1 across $[r_1,r_2]$ is
\emph{exactly}
\begin{equation}\label{eq:precfall}
  P(R_1)-P(R_2)\;\le\;g_{\mathrm{CE}}(R_2)-g_{\mathrm{CE}}(R_1)\;=:\;\text{margin},
  \tag{P}
\end{equation}
a bound on how far precautionary saving falls across the interval
(\lean{CertaintyEquivalent.saving\_le\_of\_precautionary\_fall}). No derivative appears. Two weakenings
of \eqref{eq:precfall} are worth separating, and they close at different risk aversions. Measured over
$r\in[4\%,6\%]$ at zero wealth in the top earnings state:

\begin{center}
\begin{tabular}{lrrrrrr}
\toprule
$\gamma$ & $1$ & $2$ & $2.5$ & $3$ & $3.2$ & $3.39$\\
\midrule
$P(R_1)$ & $0.109$ & $0.154$ & $0.176$ & $0.198$ & $0.206$ & $0.214$\\
$P(R_1)-P(R_2)$ & $0.044$ & $0.032$ & $0.028$ & $0.025$ & $0.024$ & $0.023$\\
margin & $0.329$ & $0.107$ & $0.062$ & $0.031$ & $0.022$ & $0.014$\\
\midrule
level route: $P(R_1)\le$ margin & $\checkmark$ & $\times$ & $\times$ & $\times$ & $\times$ & $\times$\\
$\theta$ admissible & --- & $0.305$ & $0.650$ & $0.842$ & $0.895$ & $0.936$\\
$\theta$ true & $0.598$ & $0.793$ & $0.839$ & $0.872$ & $0.883$ & $0.892$\\
\bottomrule
\end{tabular}
\end{center}

At $\gamma=1$ the crudest weakening works: discard $P(R_2)\ge0$ and a bound on the \emph{level} of
precautionary saving suffices, with the margin $0.329$ exceeding $P(R_1)=0.109$ threefold
(\lean{saving\_le\_of\_precautionary\_level}). Above that the level route dies and one needs the
multiplicative form: keeping a fraction $\theta$ of precautionary saving at the higher rate,
\eqref{eq:precfall} becomes $(1-\theta)P(R_1)\le\text{margin}$
(\lean{saving\_le\_of\_precautionary\_ratio}). At $\gamma=2$ the admissible $\theta$ is $0.305$ against
a true ratio of $0.793$: precautionary saving need only be shown not to lose seventy per cent of
itself when the rate rises by two points, which is a crude statement. At $\gamma=3$ the admissible
$\theta$ is $0.842$ against $0.872$, a margin of three per cent. Above $\gamma\approx3.1$ there is no
margin at all, because Theorem 1 itself fails on that interval --- the $3.39$ of the derivative
comparison is the threshold at $r=4\%$, not across $[4\%,6\%]$.

So the honest target is $\gamma\le2$, and what it needs is two-sided bounds on precautionary saving at
one state accurate to about a quarter: measured, $P(R_1)=0.154$ and $P(R_2)=0.122$ against a required
gap of $0.107$. The sandwich available gives $\kappa_kH_k=1.40$ at that state, nine times too loose,
so this is a new estimate rather than an assembly of existing ones --- but it is a bound on a level,
at one state, with a factor of four to give away, which is a different kind of demand from anything
the increment route leaves open. Numerics: \texttt{precbound.m}.

\subsection{The channel is the borrowing constraint}

The first-order scaling points at the constraint rather than at precaution, and the point can be
tested directly: loosen the borrowing limit and watch the threshold. If the constraint is the
channel, the risky threshold should walk back to the certainty-equivalent one as the limit
approaches the natural one, $y_{\min}/r$. At Aiyagari's calibration it does
(\texttt{numerics/borrow.m}), and it does so almost linearly:

\begingroup\small
\begin{center}
\begin{tabular}{@{}lcccccccc@{}}
\toprule
Borrowing limit & $0$ & $0.5$ & $1$ & $2$ & $3$ & $4$ & $5$ & $6$\\
\midrule
Threshold $\mu$ & $3.42$ & $3.50$ & $3.58$ & $3.74$ & $3.89$ & $4.05$ & $4.21$ & $4.38$\\
Share of the gap closed & $0\%$ & $11\%$ & $21\%$ & $40\%$ & $57\%$ & $74\%$ & $90\%$ & $104\%$\\
\bottomrule
\end{tabular}
\end{center}
\endgroup

The natural limit here is $y_{\min}/r=6.75$, and by a limit of $6$ the gap is gone: the risky
household's threshold, $4.38$, is the certainty-equivalent household's, $4.33$. So at this
calibration the borrowing constraint accounts for the whole of the correction, and precaution
accounts for none of it. Precaution is certainly present --- at the loose limit the risky
household still saves $1.62$ against the certainty-equivalent $1.46$, a precautionary excess of
eleven per cent --- but it moves the level of saving without moving the sign of its response to
the interest rate. Across the other calibrations we ran (\texttt{numerics/borrow3.m}) loosening
the limit always closes the gap and sometimes overshoots it, so the constraint is the dominant
channel rather than an exact decomposition: with ample borrowing, risk can make Theorem 1 easier
than certainty equivalence rather than harder.

For the programme this is the useful part. The missing ingredient for $\mu=3$ is not a
precautionary bound but a bound on what the borrowing constraint does to the comparative static,
and that is a more tractable object: without a binding constraint the household's Euler equation
holds with equality at every state and Theorem 1 at zero wealth reduces to the duration criterion,
which is proved outright for $\mu\le1$ and is an explicit finite computation in the primitives
otherwise. What remains is the constrained states, which at this calibration are about eight per
cent of the state space at every age.

\subsection{The constraint channel is monotone, and needs no condition}

The comparative static in the borrowing limit itself is a theorem, and unlike the one in the
interest rate it needs nothing beyond the standing assumptions.

\begin{theorem}[A tighter limit means less consumption and more saving]\label{thm:limit}
Two households share every primitive and differ only in the borrowing limit, $f_2\le f_1$. Then at
every age and every state in the tighter household's region it consumes less and saves more.
\formal{IncomeFluctuation.finiteConsumption\_le\_of\_floor\_le,
IncomeFluctuation.le\_finitePolicy\_of\_floor\_le}
\end{theorem}

The induction is the one of Theorem~\ref{thm:reduction} with the interest rate held fixed, and
that is exactly what removes the difficulty. Suppose the claim holds with $k$ periods to come and
fails at $k+1$, so the tighter household consumes strictly more and therefore saves strictly less.
Its Euler inequality under the cap, the looser household's Euler inequality above its floor ---
available because the looser household saves strictly more than the tighter one and hence strictly
more than either floor --- the induction hypothesis, and monotonicity of consumption in wealth
chain into $u'(c_1)\ge u'(c_2)$, contradicting $c_1>c_2$. The gross return appears on both sides
and cancels, so there is no substitution term to control and no analogue of
\eqref{eq:stageslack} is needed. Tightening the limit raises the marginal value of saving
unambiguously, and the terminal stage is again free: a household with a lower floor runs down
further in its last period.

That is the one-sided half of what the numbers showed, and it says the constraint channel is
better behaved than the rate channel rather than merely different. What it does not yet give is a
\emph{size}: Theorem~\ref{thm:limit} orders the two consumption functions without bounding the
distance between them, and it is the distance that turns into the $0.9$ units of risk aversion
between the constrained and certainty-equivalent thresholds. The natural next object is therefore
a modulus for Theorem~\ref{thm:limit}, and it has a clean form.

\begin{theorem}[The modulus: a relaxation of $\Delta$ is worth at most $\Delta$ of wealth]
\label{thm:modulus}
Let $\Delta=f_1-f_2\ge0$ and $r\ge0$. Then at every age and every state,
\[
  c^{f_1}_k(x)\ \le\ c^{f_2}_k(x)\ \le\ c^{f_1}_k(x+\Delta) .
\]
\formal{IncomeFluctuation.finiteConsumption\_le\_of\_floor\_le,
IncomeFluctuation.finiteConsumption\_le\_shift}
\end{theorem}

The household that may borrow down to $f_2$ therefore behaves, at wealth $x$, no more freely than
the household confined to $f_1$ would at wealth $x+\Delta$: the entire value of the extra
borrowing capacity is bounded by the capacity itself. The proof is the same backward induction,
carried out at shifted arguments, and the shift pays for itself twice. The terminal stage holds
with slack, since the looser household gains $\Delta$ by running down further while the shifted
comparison gives it $(1+r)\Delta$; and the interiority the step needs comes free, because the
failure hypothesis forces the tighter household's saving to exceed the looser one's by more than
$\Delta$, hence to exceed $f_2+\Delta=f_1$.

Two consequences. Since saving never falls with wealth, the sandwich gives the plain bound
$0\le c^{f_2}_k(x)-c^{f_1}_k(x)\le(1+r)\Delta$, so the effect of the limit on consumption is at
most the relaxation, up to one period's interest. And because the bound is stated as a shift in
wealth rather than as an additive constant, it composes with the minimal propensity to consume:
whatever upper bound on the propensity to consume out of wealth one has at age $k$, the modulus
turns it directly into a bound on the consumption gap, and it is that product which becomes the
gap between the constrained and certainty-equivalent thresholds. \subsection{The upper bound on the propensity to consume, and what it is worth}

The missing piece is an upper bound on the propensity to consume out of wealth. The repository has
the minimal propensity, which is a lower bound; the upper one follows by the same two-point Euler
comparison, with no derivative taken. Write the bound as a secant slope in assets, so that $s_k$
bounds $c_k(a_2)-c_k(a_1)$ by $s_k(a_2-a_1)$.

\begin{theorem}[The propensity to consume falls with the horizon]\label{thm:mpc}
In the last period the slope is exactly $R$ (\lean{stageConsumption\_sub\_zero}). If at stage $k$
the slope is at most $s$ and consumption is at least $c_{\min}$, then at a state where the poorer
household saves, the slope at stage $k+1$ is at most
\[
  R\cdot\frac{\nu}{1+\nu},\qquad \nu=\frac{c\,s}{c_{\min}},
\]
with $c$ the poorer household's consumption.
\formal{IncomeFluctuation.stageConsumption\_sub\_le\_of\_slope}
\end{theorem}

The step is the two Euler inequalities at the two wealth levels: the induction hypothesis bounds
the rise in tomorrow's consumption by $s\delta$ with $\delta$ the difference in saving, which
bounds the fall in marginal utility by the factor $(1+s\delta/c_{\min})^{-\mu}$, and the budget
identity $c_2-c_1=R(a_2-a_1)-\delta$ turns that into the stated slope. Since $\nu\mapsto\nu/(1+\nu)$
is a contraction towards $1-1/\nu$, the slope falls from $R$ towards $R(1-c_{\min}/(c\,s))$ as the
horizon lengthens: the longer the life, the less of an extra unit of wealth is spent at once,
which is the qualitative fact one wants.

It is not sharp enough. At Aiyagari's calibration, read at zero wealth in the highest earnings
state, consumption is about $1.32$ and the worst consumption next period about $0.88$, so
$\nu\approx1.5\,s$ and the recursion settles at a propensity of about $0.36$, against a true
propensity there of about $0.04$. Fed through Theorem~\ref{thm:modulus} it bounds the effect of
the borrowing limit on consumption by about $2.5$ where the truth is about $0.04$, a factor of
sixty. The loss is the ratio of consumption today to the worst consumption tomorrow, and that
ratio is where the sandwich's slack of order human wealth enters --- the same obstruction that
closed the infinite-horizon routes in the companion write-up, met here for the fourth time and in
the same form.

So the chain now has every link proved except its sharpness. Theorem~\ref{thm:ce} gives the
criterion without risk, Theorem~\ref{thm:limit} signs the constraint's effect,
Theorem~\ref{thm:modulus} bounds it by a shift in wealth, and Theorem~\ref{thm:mpc} converts a
shift in wealth into consumption. What none of them supplies is a bound on consumption that is
tight to a few per cent rather than to a factor of ten, and every route we have tried reduces to
exactly that.

\subsection{Carrying earnings across all remaining ages}

The two halves of the sandwich are not equally to blame. The upper half,
$c_k\le\kappa_k(m+H_k)$, carries the expected present value of future earnings and is about twenty
per cent above the truth at Aiyagari's calibration. The lower half, $\kappa_k m\le c_k$, ignores
future earnings entirely and is a factor of ten below it. The repair is the same device applied on
the other side: carry earnings forward across all the remaining ages instead of one step at a
time.

\begin{theorem}[The lower half, with human wealth]\label{thm:lower}
Let $H_k$ be the present value of $k$ periods of the minimum earnings, $H_0=0$ and
$RH_{k+1}=y_{\min}+H_k$. Then $\lambda_k(m+H_k)\le c_k(m)$ at every age and state, where
$\lambda_0=1$ and $\lambda_{k+1}=\Thorn\lambda_kR/(\Thorn+\lambda_kR)$.
\formal{IncomeFluctuation.minHumanWealth\_le\_stageConsumption}
\end{theorem}

The old bound is the case $H=0$. The propensity $\lambda_k$ is the finite-horizon propensity of
Theorem~\ref{thm:sandwich} discounted once more by patience at each step, and the extra discount is
not decoration: without it the bound is false at the borrowing constraint, where the household
consumes exactly its cash on hand while $\kappa_k(m+H_k)$ exceeds it, because $H_k$ discounts
earnings at $R$ and $1/\kappa_k$ discounts at $R/\Thorn$. Discounting $H_k$ at $R/\Thorn$ instead
repairs the corner and breaks the induction; the extra patience factor repairs both at once
(\lean{lowMPC\_mul\_minHumanWealth\_le}), and at $\beta=0.96$, $r=4\%$ and $\mu=3$ it costs about
three per cent over sixty periods.

At the state where the life cycle binds the household enters the next period with cash on hand
about $2.0$ and consumes about $0.88$. The old bound gives $0.086$; this one gives $0.33$. Through
Theorem~\ref{thm:mpc} that improves the bound on the propensity to consume from about $0.94$ to
about $0.75$, against a true propensity of $0.04$.

So the multi-age argument buys a factor of four and not a factor of ten, and the reason is visible
in the statement: it carries the \emph{minimum} earnings, and a household that drew the minimum
forever really would consume this little. What the programme needs is a lower bound carrying the
\emph{expected} human wealth, and that is not available from any chain of Euler inequalities,
because the Euler inequality at a constrained state is exactly the statement that expected future
earnings cannot be borrowed against. That is the one quantitative fact about a single household
on which the whole programme turns, and identifying it precisely is what makes it attackable.

\subsection{The constraint-incidence bound}

Both halves can be carried at once. Accumulate human wealth by the \emph{risk-adjusted}
expectation --- the power mean with exponent $-\mu$, which is what the Euler equation itself
produces --- and then \emph{cap} it state by state at the level the borrowing constraint permits:
\[
  H_0=0,\qquad
  H_{k+1}(z)=\min\Bigl\{\tfrac1R\,M^z_{-\mu}\bigl(y_w+H_k(w)\bigr),\
  \Bigl(\tfrac1{\kappa_{k+1}}-1\Bigr)y_z\Bigr\}
\]
(\lean{riskHumanWealth}), with $\kappa$ the propensity of Theorem~\ref{thm:sandwich}.

\begin{theorem}[The constraint-incidence bound]\label{thm:incidence}
$\kappa_k\bigl(m+H_k(z)\bigr)\le c_k(a,z)$ at every age and state.
\formal{IncomeFluctuation.riskHumanWealth\_le\_stageConsumption}
\end{theorem}

The two ingredients meet exactly where they should. The power mean is what the Euler inequality
delivers, and the step needs only that it is superadditive under a shift, which follows from its
concavity and homogeneity (\lean{powerMean\_add\_const\_le}). The cap is the constraint's
incidence: at a corner the household consumes its cash on hand, which is at least $y_z$, and the
cap is precisely the largest human wealth the bound may claim without exceeding that. No extra
patience factor is needed, because the cap does the work the patience factor did in
Theorem~\ref{thm:lower}.

At Aiyagari's calibration, in the first period of a sixty-period life
(\texttt{numerics/incidence.m}), the cap binds at the two lowest earnings states and nowhere else,
which is where a household at zero assets is against the limit. The bound reaches

\begingroup\small
\begin{center}
\begin{tabular}{@{}lccc@{}}
\toprule
At zero wealth & lowest earnings & middle & highest\\
\midrule
Bound as a fraction of the truth & $1.000$ & $0.890$ & $0.877$\\
Minimum-earnings bound, for comparison & $1.010$ & $0.409$ & $0.295$\\
\bottomrule
\end{tabular}
\end{center}
\endgroup

and over the whole grid it never exceeds the truth, touching it exactly at the constrained corner.
So the consumption bound is now within about twelve per cent where the life cycle binds, against a
factor of three before. That is the sharpest bound in the development, and it was reached by
bounding where the constraint binds rather than by another Euler chain.

It does not improve the propensity bound of Theorem~\ref{thm:mpc}, because that bound is driven by
the worst earnings state next period, which is exactly where the cap binds. Closing the threshold
argument would need the same sharpening on the worst state, and that turns out to be impossible
within this class.

\subsection{The worst state is pinned}

Two facts settle it. The first is that the cap is not conservatism but necessity.

\begin{theorem}[The cap is necessary]\label{thm:cap}
Wherever a household at zero assets is against the borrowing limit, every affine lower bound
$\kappa_k(m+H)$ valid at that state obeys $H\le(1/\kappa_k-1)y_z$.
\formal{IncomeFluctuation.cap\_necessary\_of\_corner}
\end{theorem}

The proof is one line --- at the corner consumption is exactly $y_z$ --- and the content is that
Theorem~\ref{thm:incidence} claims the largest human wealth any bound of its shape may claim. It
cannot be sharpened within the affine class at all.

The second is that dropping the cap does not merely weaken the proof, it makes the bound false.
Measuring effective human wealth as $c/\kappa-m$ at next-period assets of $1.67$
(\texttt{numerics/worst.m}):

\begingroup\small
\begin{center}
\begin{tabular}{@{}lcccc@{}}
\toprule
Effective human wealth at & truth & capped mean & uncapped mean & expectation\\
\midrule
lowest earnings & $10.96$ & $5.99$ & $13.11$ & $16.57$\\
highest earnings & $27.92$ & $23.90$ & $26.38$ & $32.24$\\
\bottomrule
\end{tabular}
\end{center}
\endgroup

At the lowest earnings state the uncapped power mean, $13.11$, \emph{exceeds} the truth, $10.96$,
so a bound built on it is invalid: the household will be against the limit later, and the uncapped
recursion cannot see that. The truth sits between the two, and the distance is the price of
capping at zero assets when the household in fact holds some. A bound that closed it would have to
let human wealth depend on the assets held, which requires a lower bound on the household's own
asset path; the one available, from the upper sandwich, is negative at the state that matters.

So the worst state is pinned from both sides: the cap cannot be raised, because the corner forbids
it, and it cannot be dropped, because the truth is below what dropping it would claim. That closes
this family of arguments, and with it, for now, the route from a consumption bound to the
threshold. What the programme has instead is an exact statement of what a sharper argument must
supply: a lower bound on consumption that is asset-dependent rather than affine, and the asset
path to support it.

\subsection{The bound with assets in it}\label{sec:assetdep}

The obstruction just described is no longer binding, because the path bound it asked for is
available. Proposition~\ref{prop:floor} read at the household's own earnings state rather than at
the worst one is positive exactly where the cap bites hardest: a household in a high earnings
state saves, so a period later it is not at the corner, and the cap that then applies to it is the
one for a household with assets.

So assets go through the recursion twice over. Writing $\Phi_k(z,a)$ for that saving floor,
clamped at zero,
\[
  G_0(z,a)=0,\qquad
  G_{k+1}(z,a)=\min\Bigl(\tfrac1R M^z_{-\mu}\bigl(y_w+G_k(w,\Phi_{k+1}(z,a))\bigr),\;
    \bigl(\tfrac1{\kappa_{k+1}}-1\bigr)(y_z+Ra)\Bigr),
\]
and the bound is $\kappa_k(m+G_k(z,a))\le c_k(a,z)$. The cap now grows with assets, and tomorrow's
human wealth is read at the assets the household will hold rather than at zero.

\begin{theorem}[The asset-dependent incidence bound]\label{thm:assetdep}
$\kappa_k\bigl(m+G_k(z,a)\bigr)\le c_k(a,z)$, and $G_k(z,a)\ge H_k(z)$ for $a\ge0$, so the bound is
never worse than Theorem~\ref{thm:incidence}.
\formal{IncomeFluctuation.assetHumanWealth\_le\_stageConsumption},
\lean{IncomeFluctuation.riskHumanWealth\_le\_assetHumanWealth}
\end{theorem}

Measured against the truth over the whole grid at three ages
(\texttt{numerics/assetdep.m}, \texttt{numerics/worstgrid.m}):

\begingroup\small
\begin{center}
\begin{tabular}{@{}lccc@{}}
\toprule
Worst ratio of bound to truth & $r=-7\%$ & $r=4\%$ & $r=10\%$\\
\midrule
Theorem~\ref{thm:incidence}, cap at zero assets & $0.237$ & $0.606$ & $0.760$\\
Theorem~\ref{thm:assetdep}, cap at the assets held & $0.350$ & $0.746$ & $0.857$\\
\bottomrule
\end{tabular}
\end{center}
\endgroup

The new bound never exceeds the truth anywhere on the grid. The gain is largest exactly where the
old one was weakest, at positive assets in low earnings states: at $r=4\%$ with five units of
assets and the lowest earnings, the ratio rises from $0.621$ to $0.835$. Where the old bound was
already good, at zero assets, little changes, which is as it should be, since that is the state
the cap was computed for.

What the refinement does not improve is the ceiling on aggregate capital, because the state that
binds there is the highest earnings one, where the cap was never active: the implied ceiling at
$r=-7\%$ is $22.15$ against $22.17$. That is still far below the $39.66$ that feasibility alone
gives, which is worth carrying into the existence argument of Section~\ref{sec:eqm}.

\section{Theorem 1 at logarithmic utility}\label{sec:log}

The slack condition of Section~\ref{sec:eqm} is an Euler condition, and it is worth asking whether
it settles the one case the stationary model does settle, $\mu=1$. It does not, and the reason is
instructive. Checking it on the whole grid at $\gamma=1$ (\texttt{numerics/slacklog.m}), it holds
at $r=0$ and $r=2\%$ and fails at $r=4\%$ and above, by about twelve per cent in elasticity terms,
at the youngest age with zero assets in the highest earnings state. Since a life-cycle equilibrium
sits above the rate of time preference, $4.17\%$ here, the Euler reduction says nothing where it
is wanted, even at logarithmic utility.

\subsection{A condition on the value gap}

A different sufficient condition does hold. Write $V_k(\cdot;R)$ for the stage-$k$ value.

\begin{definition}[Value gap]\label{def:gap}
The value gap has slope at least $\theta$ at stage $k$ if, for $x\le y$ in $[0,\bar a]$,
\[
  \theta\,(y-x)\ \le\ \bigl[V_k(y;R_2)-V_k(y;R_1)\bigr]-\bigl[V_k(x;R_2)-V_k(x;R_1)\bigr].
\]
\formal{IncomeFluctuation.ValueGap}
\end{definition}

By the envelope theorem the slope of the gap is $R_2u'(c_2)-R_1u'(c_1)$, so the condition says the
marginal value of assets rises with the rate. At $\gamma=1$ that reads $c(R_2)/R_2\le c(R_1)/R_1$:
consumption rises by proportionally less than the rate does.

\begin{theorem}[Theorem 1 from the value gap]\label{thm:gapmono}
If the value gap has a positive slope at stage $k$, then saving at stage $k$ rises with the rate,
at every state. Aggregate capital rises with it.
\formal{IncomeFluctuation.stagePolicy\_mono\_of\_valueGap},
\lean{IncomeFluctuation.olgCapital\_le\_of\_valueGap}
\end{theorem}

Unlike the Euler route this needs no induction over ages. The condition on the continuation gives
the conclusion for the household that faces it, directly, and the argument is Topkis with two
continuations: the two households differ only in their cash on hand, which is ordered by the rate,
and in their continuations, whose difference is increasing, so the higher-rate household cannot
save less. Only concavity of period utility is used, as in Section~\ref{sec:ret}.

\subsection{Why it holds at $\mu=1$, and only near there}

At $\gamma=1$ the patience factor is $\text{\th}=\beta R$, so $q=\text{\th}/R=\beta$ and the
propensity $\kappa_k=1/\sum_{i\le k}\beta^i$ of Proposition~\ref{prop:kappa} does not depend on
the interest rate at all. Human wealth does, and falls as the rate rises. Where the borrowing
constraint is slack, $c=\kappa_k(\bar y+Ra+H)$ and
\[
  R_2c(R_1)-R_1c(R_2)\ =\ \kappa_k\bigl[(R_2-R_1)\bar y+R_2H_1-R_1H_2\bigr]\ >\ 0 ,
\]
which is the condition, with room. At a corner, $c=\bar y+Ra$ and the same inequality reduces to
$R_1\bar y\le R_2\bar y$.

Numerically (\texttt{numerics/condB.m}) the condition holds at $\gamma=1$ at every age, state and
rate tested, and fails from $\gamma=2$, where between $19\%$ and $51\%$ of states violate it. It is
therefore sufficient and not necessary: Theorem 1 itself still holds at $\gamma=2$, and up to about
$3.4$. So this route covers $\mu=1$ and closes nothing above it.

What is proved, then, is the reduction, and it is a reduction to a condition that is true at
logarithmic utility rather than one that is false there. Discharging the condition itself from
primitives is not available from the sandwich of Theorem~\ref{thm:sandwich}: the sandwich's width
is $\kappa_kH_k$, and the rate response of exactly that quantity is what drives the gap, so the
bounds cannot see it. That is the same obstruction the consumption-bound family ran into above,
and it is the piece a sharper argument would have to supply.

\section{What remains}\label{sec:prog}

\subsection{Above logarithmic utility}

The one thing the programme set out to do and has not done is uniqueness at $\mu\in\{3,5\}$. The
argument of Section~\ref{sec:inc} does not fail there for want of propagation: the per-stage
multiplier $m_k=qA/(\Thorn+qA\kappa_kR)$ has a product below one over a life at every risk aversion we
tested, up to $\gamma=5$. What gives way is the certainty-equivalent side. The gradient sum $A$ is
one plus the variance of the risk tilt at $\gamma=1$ and only bounded by
$(1+\operatorname{Var})^{(\gamma+1)/(2\gamma)}$ above it, by Lyapunov, and both the loss in that
inequality and the variance itself grow with $\gamma$ --- the variance tenfold between $\gamma=1$ and
$\gamma=3$. The bound crosses the allowance between $\gamma=1$ and $\gamma=1.5$.

So the question at $\mu=3$ is whether $\mathbb E_w[u^{1+1/\gamma}]$ can be bounded better than by
Lyapunov, given $\mathbb E_w[u]=1$ and a bound on the variance. That is a question about moments of a
positive random variable with known mean, not about households, and it is where we would look next.

\subsection{Routes that are closed}

Three, each with a reason worth keeping so that none is reopened.

\paragraph{Mean field games.} Section~\ref{sec:mfg}. Date-by-date Lasry--Lions monotonicity is false
here at every risk aversion, and the summed condition that the argument actually uses, taken to its
weakest form, collapses to monotone capital supply. What survives is about the life-cycle household:
the reduction of the summed pairing to a mixed difference, the duration wedge for $\Phi$, and the
exact re-pricing identity.

\paragraph{Stochastic dominance.} Light and Weintraub's Theorem 2 reaches uniqueness from ordered
distributions, and the obvious weakening of its hypothesis is to ask only for that ordering: first-order
stochastic dominance of the cohort distributions. It buys nothing. In wage units, at $r=4\%$ against
$6\%$, dominance holds to machine precision at $\gamma\le3$, where the policy ordering already holds
exactly, and fails at precisely the risk aversions where the policy ordering fails --- at $\gamma=5$,
eleven of sixty ages have a crossing of the asset CDF, worst gap $5\times10^{-2}$, and $0.57$
conditional on the earnings state. The distributions genuinely cross, and they stop being ordered
exactly when the policies do, while aggregate supply stays monotone throughout. The aggregate is the
weaker statement and there is no distributional route to it (\texttt{fosdcheck.m}).

\paragraph{Two-sided sandwiches.} For the two-rate comparison at a state, the width of the sandwich at
a given rate \emph{is} the precautionary wedge: arithmetic human wealth gives the ceiling,
risk-adjusted gives the floor, and they cannot both be tight. So bounding the two problems separately
works only while the consumption response to the rate exceeds that wedge, which at Aiyagari's numbers
is $\gamma=1$ alone (Section~\ref{sec:thm1}). Any route to $\mu>1$ has to bound the \emph{difference}
between the two problems, which is what Section~\ref{sec:inc} does.

\subsection{The classical two-period results}

Section~\ref{sec:diamond} does the steady state, and for every risk aversion rather than only for
logarithmic utility. What remains of the two-period economy is standalone: monotone global
convergence, existence under Inada conditions with a multiplicity witness in general, dynamic
inefficiency of a steady state with $r<n$ by an explicit Pareto-improving transfer, the golden rule,
and the autarkic and monetary steady states of the pure-exchange economy.

\subsection{Extensions}

An age profile of earnings; a government budget alongside the pension of Section~\ref{sec:ret};
survival risk with accidental bequests as a second fixed point; population growth. Each of these
changes the equilibrium object rather than the household, so the existence argument of
Section~\ref{sec:eqm} should carry with the bracket recomputed.

\subsection{The certificate}

The uniqueness result of Section~\ref{sec:inc} rests on a hypothesis discharged by a computation in
exact rational arithmetic rather than by a Lean proof term. Two things would improve its standing. The
computation could be replayed inside Lean, which for a branch and bound over rational boxes is a
matter of \texttt{Decidable} arithmetic and a large \texttt{native\_decide}, or the variance bound could
be replaced by an analytic argument, which is the same question as the one above about moments. Until
one of those, the result has the standing of the existence bracket of Section~\ref{sec:eqm}: checked,
and not checked by the kernel.


\bibliographystyle{ecta}
\bibliography{refs}

\appendix
\section{Index of formal results}\label{app:index}

All in the namespace \lean{LeanEconomics.IncomeFluctuation}, directory \lean{LeanEconomics/OLG}.
Every result depends only on the three standard axioms (\lean{propext}, \lean{Classical.choice},
\lean{Quot.sound}) and none uses \lean{sorry}.

\begingroup\footnotesize
\begin{longtable}{@{}p{0.38\textwidth}p{0.40\textwidth}p{0.16\textwidth}@{}}
\toprule
Result & Lean name & Module \\
\midrule
\endhead
Stage value, policy, consumption & \lean{stageValue}, \lean{stagePolicy}, \lean{stageConsumption} & OLG/\allowbreak LifeCycle\\
Last period: no saving & \lean{stagePolicy\_zero} & OLG/\allowbreak LifeCycle\\
Last period: consume cash on hand & \lean{stageConsumption\_zero} & OLG/\allowbreak LifeCycle\\
Positive consumption at every stage & \lean{stageConsumption\_pos} & OLG/\allowbreak LifeCycle\\
Saving and consumption rise with wealth & \lean{stagePolicy\_mono}, \lean{stageConsumption\_mono} & OLG/\allowbreak LifeCycle\\
Carroll--Kimball at every stage & \lean{concaveOn\_stageConsumption} & OLG/\allowbreak LifeCycle\\
Euler under the cap, by stage & \lean{stage\_euler\_le} & OLG/\allowbreak LifeCycle\\
Euler above the floor, by stage & \lean{stage\_euler\_ge} & OLG/\allowbreak LifeCycle\\
Finite-horizon propensities, human wealth & \lean{stageMPC}, \lean{stageHumanWealth} & OLG/\allowbreak LifeCycle\\
Propensities in $(0,1]$, above $\kappa$ & \lean{stageMPC\_pos}, \lean{stageMPC\_le\_one}, \lean{minMPC\_le\_stageMPC} & OLG/\allowbreak LifeCycle\\
Lower Euler step & \lean{crra\_stageMPC\_step} & OLG/\allowbreak LifeCycle\\
Upper Euler step & \lean{crra\_stageUpper\_step} & OLG/\allowbreak LifeCycle\\
Theorem~\ref{thm:sandwich}, lower & \lean{stageMPC\_mul\_le\_stageConsumption} & OLG/\allowbreak LifeCycle\\
Theorem~\ref{thm:sandwich}, upper & \lean{stageConsumption\_le\_stageHumanWealth} & OLG/\allowbreak LifeCycle\\
One-step saving bound & \lean{stagePolicy\_le} & OLG/\allowbreak LifeCycle\\
\midrule
Retired household (lump-sum pension) & \lean{withPension} & OLG/\allowbreak Retirement\\
Cash on hand rises with assets & \lean{withPension\_resources\_strictMono} & OLG/\allowbreak Retirement\\
Human wealth is state-independent & \lean{stageHumanWealth\_withPension\_const} & OLG/\allowbreak Retirement\\
Perpetuity bound $H_k\le p/r$ & \lean{stageHumanWealth\_withPension\_le} & OLG/\allowbreak Retirement\\
Life-cycle value with retirement & \lean{lifeValue} & OLG/\allowbreak Retirement\\
Saving, consumption rise with wealth & \lean{lifePolicy\_mono}, \lean{lifeConsumption\_mono} & OLG/\allowbreak Retirement\\
\midrule
Proposition~\ref{prop:tangent} & \lean{crra\_le\_tangent} & OLG/\allowbreak Diamond\\
Closed-form saving & \lean{saveShare}, \lean{saveShare\_euler} & OLG/\allowbreak Diamond\\
It is the optimum & \lean{saveShare\_isOptimal} & OLG/\allowbreak Diamond\\
Proposition~\ref{prop:diamondcs} & \lean{saveShare\_strictMonoOn}, \lean{saveShare\_strictAntiOn} & OLG/\allowbreak Diamond\\
Proposition~\ref{prop:engine} & \lean{div\_one\_add\_rpow\_strictMono} & OLG/\allowbreak Diamond\\
Theorem~\ref{thm:diamond} & \lean{exists\_unique\_steadyState} & OLG/\allowbreak Diamond\\
\midrule
Theorem~\ref{thm:detclosed} & \lean{stageConsumption\_eq\_of\_interior} & OLG/\allowbreak Deterministic\\
Rate-free propensity at log & \lean{stageMPC\_one\_eq\_of\_rate} & OLG/\allowbreak Deterministic\\
Human wealth falls with the rate & \lean{stageHumanWealth\_antitone\_rate} & OLG/\allowbreak Deterministic\\
Theorem~\ref{thm:detlog} & \lean{stagePolicy\_mono\_det\_log}, \lean{olgCapital\_mono\_det\_log} & OLG/\allowbreak Deterministic\\
\midrule
Chord of a concave power & \lean{CertaintyEquivalent.chord\_le\_rpow} & OLG/\allowbreak Certainty\allowbreak Equivalent\\
Theorem~\ref{thm:chord} & \lean{CertaintyEquivalent.consumption\_antitone\_of\_chord} & OLG/\allowbreak Certainty\allowbreak Equivalent\\
\midrule
Pension schedules & \lean{lumpSum}, \lean{tapered}, \lean{cliff} & OLG/\allowbreak MeansTest\\
Cash on hand & \lean{cash} & OLG/\allowbreak MeansTest\\
Theorem~\ref{thm:topkis} & \lean{isOptimalSaving\_mono} & OLG/\allowbreak MeansTest\\
Monotone in assets when cash is & \lean{isOptimalSaving\_mono\_assets} & OLG/\allowbreak MeansTest\\
Gentle taper keeps cash rising & \lean{cash\_strictMono\_tapered} & OLG/\allowbreak MeansTest\\
A cliff reverses cash on hand & \lean{cash\_lt\_cash\_of\_cliff} & OLG/\allowbreak MeansTest\\
Theorem~\ref{thm:cons} & \lean{exists\_consumption\_decrease} & OLG/\allowbreak MeansTest\\
Theorem~\ref{thm:taper}, cliff & \lean{exists\_saving\_decrease\_in\_assets} & OLG/\allowbreak MeansTest\\
Theorem~\ref{thm:azmono} & \lean{policyOf\_mono\_of\_lt} & OLG/\allowbreak MeansTest\\
\midrule
One age's expectation operator & \lean{stageStep} & OLG/\allowbreak Equilibrium\\
Monotone class preserved & \lean{monoRegion\_stageStep} & OLG/\allowbreak Equilibrium\\
Ordered policies, ordered functions & \lean{stageStep\_le\_stageStep} & OLG/\allowbreak Equilibrium\\
Life-cycle assets, aggregate capital & \lean{cohortAssets}, \lean{olgCapital} & OLG/\allowbreak Equilibrium\\
Theorem~\ref{thm:thm2} & \lean{cohortAssets\_le}, \lean{olgCapital\_le\_of\_stagePolicy\_le} & OLG/\allowbreak Equilibrium\\
Proposition~\ref{prop:floor} & \lean{le\_stagePolicy} & OLG/\allowbreak Equilibrium\\
Theorem~\ref{thm:thm3} & \lean{olgEquilibriumRate\_unique} & OLG/\allowbreak Equilibrium\\
\midrule
Argmax at a concave continuation & \lean{argmax\_eq\_singleton\_of\_concaveSlices} & OLG/\allowbreak Continuity\\
Slice commutes with the operator & \lean{sliceAt\_bellman}, \lean{sliceAt\_augStageValue} & OLG/\allowbreak Continuity\\
Augmented stage policy & \lean{augStagePolicy}, \lean{augStagePolicy\_eq} & OLG/\allowbreak Continuity\\
Jointly continuous stage policy & \lean{continuousOn\_augStagePolicy} & OLG/\allowbreak Continuity\\
Cohort recursion in the augmented state & \lean{augCohortAssets}, \lean{augCohortAssets\_eq} & OLG/\allowbreak Continuity\\
Theorem~\ref{thm:cont} & \lean{continuousOn\_olgSupply} & OLG/\allowbreak Continuity\\
Theorem~\ref{thm:exist} & \lean{exists\_olgEquilibrium} & OLG/\allowbreak Continuity\\
\midrule
Saving rises iff consumption responds little & \lean{stagePolicy\_le\_iff\_consumption\_le} & OLG/\allowbreak RateMonotone\\
Terminal stage, unconditional & \lean{stagePolicy\_mono\_zero} & OLG/\allowbreak RateMonotone\\
Condition \eqref{eq:stageslack} & \lean{StageSlack} & OLG/\allowbreak RateMonotone\\
Theorem~\ref{thm:reduction}, step & \lean{stagePolicy\_mono\_step} & OLG/\allowbreak RateMonotone\\
Theorem~\ref{thm:reduction}, induction & \lean{stagePolicy\_mono\_of\_stageSlack} & OLG/\allowbreak RateMonotone\\
Theorem~\ref{thm:reduction}, aggregate & \lean{olgCapital\_le\_of\_stageSlack} & OLG/\allowbreak RateMonotone\\
\midrule
Human wealth, price of the path & \lean{humanWealth}, \lean{pathPrice} & OLG/\allowbreak Certainty\allowbreak Equivalent\\
Closed form $c_0=W/D$ & \lean{consumption\_mul\_pathPrice} & OLG/\allowbreak Certainty\allowbreak Equivalent\\
Saving rises iff the cross product falls & \lean{consumption\_antitone\_iff} & OLG/\allowbreak Certainty\allowbreak Equivalent\\
Theorem~\ref{thm:ce} & \lean{consumption\_antitone\_of\_le\_one} & OLG/\allowbreak Certainty\allowbreak Equivalent\\
\midrule
Same household, different limit & \lean{withFloor} & OLG/\allowbreak Borrowing\allowbreak Limit\\
Stages at a general limit & \lean{finiteValue}, \lean{finitePolicy}, \lean{finiteConsumption} & OLG/\allowbreak Borrowing\allowbreak Limit\\
Theorem~\ref{thm:limit}, consumption & \lean{finiteConsumption\_le\_of\_floor\_le} & OLG/\allowbreak Borrowing\allowbreak Limit\\
Theorem~\ref{thm:limit}, saving & \lean{le\_finitePolicy\_of\_floor\_le} & OLG/\allowbreak Borrowing\allowbreak Limit\\
Theorem~\ref{thm:modulus} & \lean{finiteConsumption\_le\_shift} & OLG/\allowbreak Borrowing\allowbreak Limit\\
\midrule
Terminal slope is $R$ & \lean{stageConsumption\_sub\_zero} & OLG/\allowbreak MPCBound\\
Theorem~\ref{thm:mpc} & \lean{stageConsumption\_sub\_le\_of\_slope} & OLG/\allowbreak MPCBound\\
\midrule
Guaranteed human wealth & \lean{minHumanWealth}, \lean{lowMPC} & OLG/\allowbreak Lower\allowbreak Sandwich\\
The corner is respected & \lean{lowMPC\_mul\_minHumanWealth\_le} & OLG/\allowbreak Lower\allowbreak Sandwich\\
Theorem~\ref{thm:lower} & \lean{minHumanWealth\_le\_stageConsumption} & OLG/\allowbreak Lower\allowbreak Sandwich\\
\midrule
Power mean under a shift & \lean{powerMean\_add\_const\_le} & OLG/\allowbreak Incidence\\
Risk-adjusted, capped human wealth & \lean{riskHumanWealth} & OLG/\allowbreak Incidence\\
Theorem~\ref{thm:incidence} & \lean{riskHumanWealth\_le\_stageConsumption} & OLG/\allowbreak Incidence\\
Theorem~\ref{thm:cap} & \lean{cap\_necessary\_of\_corner} & OLG/\allowbreak Incidence\\
\midrule
Definition~\ref{def:regions} & \lean{StageRegions}, \lean{allRegions} & OLG/\allowbreak Life\allowbreak Cycle\\
Sandwich on a reachable family & \lean{stageConsumption\_le\_stageHumanWealth\_on} & OLG/\allowbreak Life\allowbreak Cycle\\
Floor on a reachable family & \lean{le\_stagePolicy\_on} & OLG/\allowbreak Equilibrium\\
Crude reachable bound & \lean{reach}, \lean{reach\_mono} & OLG/\allowbreak Reachable\\
Proposition~\ref{prop:reach} & \lean{reachRegions} & OLG/\allowbreak Reachable\\
Sandwich along a cohort's path & \lean{stageConsumption\_le\_stageHumanWealth\_reach} & OLG/\allowbreak Reachable\\
Floor along a cohort's path & \lean{le\_stagePolicy\_reach} & OLG/\allowbreak Equilibrium\\
Floor at the household's own state & \lean{le\_stagePolicy\_state\_on} & OLG/\allowbreak Equilibrium\\
\midrule
Proposition~\ref{prop:ceiling} & \lean{olgCapital\_le\_of\_invariant} & OLG/\allowbreak Existence\\
Affine coefficients & \lean{cohortSlope}, \lean{cohortFloor} & OLG/\allowbreak Existence\\
Proposition~\ref{prop:affine} & \lean{cohortFloor\_add\_mul\_le\_cohortAssets}, \lean{le\_olgCapital} & OLG/\allowbreak Existence\\
Mean human wealth recursion & \lean{sum\_mul\_stageHumanWealth\_succ} & OLG/\allowbreak Existence\\
Theorem~\ref{thm:patient} & \lean{stageMPC\_mul\_sum\_stageHumanWealth\_le} & OLG/\allowbreak Existence\\
The aggregate floor climbs & \lean{sum\_mul\_cohortFloor\_mono}, \lean{sum\_mul\_cohortFloor\_nonneg} & OLG/\allowbreak Existence\\
Theorem~\ref{thm:exist2} & \lean{exists\_olgEquilibrium\_of\_bounds}, \lean{exists\_olgEquilibrium\_rate} & OLG/\allowbreak Existence\\
\midrule
Proposition~\ref{prop:kappa} & \lean{stageMPC\_mul\_mpcSum}, \lean{stageMPC\_mul\_one\_sub\_pow} & OLG/\allowbreak Life\allowbreak Cycle\\
The geometric sum is bounded & \lean{sub\_mul\_mpcSum\_le\_one} & OLG/\allowbreak Life\allowbreak Cycle\\
Constant-coefficient slope & \lean{cohortSlope\_mul\_mpcSum\_succ} & OLG/\allowbreak Existence\\
Proposition~\ref{prop:slope} & \lean{mul\_geomSum\_le\_cohortSlope} & OLG/\allowbreak Existence\\
Human wealth in closed form & \lean{sum\_mul\_stageHumanWealth\_eq} & OLG/\allowbreak Existence\\
The floor in closed form & \lean{geomFloor\_le\_sum\_mul\_cohortFloor} & OLG/\allowbreak Existence\\
Theorem~\ref{thm:exist3} & \lean{exists\_olgEquilibrium\_of\_geomBounds} & OLG/\allowbreak Existence\\
\midrule
Definition~\ref{def:gap} & \lean{ValueGap} & OLG/\allowbreak Rate\allowbreak Monotone\\
Theorem~\ref{thm:gapmono} & \lean{stagePolicy\_mono\_of\_valueGap}, \lean{stagePolicy\_mono\_of\_valueGap\_all} & OLG/\allowbreak Rate\allowbreak Monotone\\
Capital from the value gap & \lean{olgCapital\_le\_of\_valueGap} & OLG/\allowbreak Rate\allowbreak Monotone\\
\midrule
The guaranteed asset path & \lean{assetFloor}, \lean{assetFloor\_le\_stagePolicy} & OLG/\allowbreak Asset\allowbreak Incidence\\
Human wealth with assets in it & \lean{assetHumanWealth}, \lean{assetHumanWealth\_mono} & OLG/\allowbreak Asset\allowbreak Incidence\\
Theorem~\ref{thm:assetdep} & \lean{assetHumanWealth\_le\_stageConsumption} & OLG/\allowbreak Asset\allowbreak Incidence\\
It dominates the asset-free bound & \lean{riskHumanWealth\_le\_assetHumanWealth} & OLG/\allowbreak Asset\allowbreak Incidence\\
\midrule
Proposition~\ref{prop:ceiling2} & \lean{cohortAssets\_le\_of\_incidence}, \lean{olgCapital\_le\_of\_incidence} & OLG/\allowbreak Existence\\
Existence with the sharper ceiling & \lean{exists\_olgEquilibrium\_of\_sharpBounds} & OLG/\allowbreak Existence\\
Proposition~\ref{prop:affceil} & \lean{cohortCeil}, \lean{cohortAssets\_le\_cohortCeil} & OLG/\allowbreak Existence\\
Ceiling on aggregate capital & \lean{olgCapital\_le\_cohortCeil} & OLG/\allowbreak Existence\\
Theorem~\ref{thm:exist4} & \lean{exists\_olgEquilibrium\_affine} & OLG/\allowbreak Existence\\
Proposition~\ref{prop:clampfloor} & \lean{cohortFloorAt}, \lean{cohortFloorAt\_le\_cohortAssets} & OLG/\allowbreak Existence\\
Floor on capital, clamped & \lean{le\_olgCapital\_at} & OLG/\allowbreak Existence\\
Theorem~\ref{thm:exist5} & \lean{exists\_olgEquilibrium\_clamped} & OLG/\allowbreak Existence\\
\midrule
The summed pairing & \lean{summedPairing}, \lean{SummedLasryLionsMonotone} & Equilibrium/\allowbreak Lasry\allowbreak Lions\\
Theorem~\ref{thm:mfgengine} & \lean{summedPairing\_nonpos\_of\_pathEquilibrium} & Equilibrium/\allowbreak Lasry\allowbreak Lions\\
Date-by-date is a special case & \lean{summedLasryLionsMonotone\_of\_forall}, \lean{summedPairing\_eq\_zero\_of\_pathEquilibrium} & Equilibrium/\allowbreak Lasry\allowbreak Lions\\
Theorem~\ref{thm:mfgagg} & \lean{aggregate\_eq\_of\_pathEquilibrium} & Equilibrium/\allowbreak Lasry\allowbreak Lions\\
Costs seeing the whole family & \lean{famPathCost}, \lean{IsFamPathEquilibrium}, \lean{famSummedPairing} & Equilibrium/\allowbreak Lasry\allowbreak Lions\\
The engine in that generality & \lean{famSummedPairing\_nonpos\_of\_pathEquilibrium} & Equilibrium/\allowbreak Lasry\allowbreak Lions\\
Theorem~\ref{thm:mfgred} & \lean{famLifetime}, \lean{famSummedPairing\_scalar}, \lean{statistic\_eq\_of\_increasingDifferences} & Equilibrium/\allowbreak Lasry\allowbreak Lions\\
Proposition~\ref{prop:loss} & \lean{famCost}, \lean{famSummedPairing\_eq\_neg\_losses} & Equilibrium/\allowbreak Lasry\allowbreak Lions\\
Each cross-loss is nonnegative & \lean{famCost\_le\_of\_pathEquilibrium}, \lean{famSummedPairing\_nonpos\_of\_losses} & Equilibrium/\allowbreak Lasry\allowbreak Lions\\
Proposition~\ref{prop:collapse} & \lean{monotone\_of\_comonotone\_of\_strictMono} & Equilibrium/\allowbreak Lasry\allowbreak Lions\\
\midrule
The accumulated gap & \lean{gapAccum}, \lean{gapAccum\_nonneg} & OLG/\allowbreak Increment\\
Theorem~\ref{thm:envelope} & \lean{negPart\_le\_gapAccum}, \lean{neg\_gapAccum\_le} & OLG/\allowbreak Increment\\
Aggregate form, single crossing & \lean{sum\_neg\_gapAccum\_le}, \lean{lt\_of\_increment\_bound} & OLG/\allowbreak Increment\\
\midrule
Superadditive certainty equivalent & \lean{powerMean\_add\_const\_le} & OLG/\allowbreak Incidence\\
Theorem~\ref{thm:mpcfloor} & \lean{stageMPC\_mul\_le\_stageConsumption\_sub} & OLG/\allowbreak Propensity\allowbreak Floor\\
Theorem~\ref{thm:cohortlip} & \lean{cohortAssets\_sub\_le\_cohortSlope} & OLG/\allowbreak Propensity\allowbreak Floor\\
Theorem~\ref{thm:marginal} & \lean{incomeStep}, \lean{stageStep\_of\_incomeFn}, \lean{sum\_mul\_incomeStep} & OLG/\allowbreak Mass\allowbreak Weighted\\
The marginal at every age & \lean{sum\_mul\_incomeStep\_iterate}, \lean{sum\_mul\_le\_of\_le} & OLG/\allowbreak Mass\allowbreak Weighted\\
Theorem 1 from the CE criterion & \lean{saving\_nonneg\_of\_correction\_le} & OLG/\allowbreak Certainty\allowbreak Equivalent\\
The precautionary split & \lean{saving\_le\_of\_precautionary\_fall} & OLG/\allowbreak Certainty\allowbreak Equivalent\\
Level and ratio routes & \lean{saving\_le\_of\_precautionary\_level}, \lean{saving\_le\_of\_precautionary\_ratio} & OLG/\allowbreak Certainty\allowbreak Equivalent\\
Level route from a floor & \lean{precautionary\_level\_of\_floor} & OLG/\allowbreak Certainty\allowbreak Equivalent\\
Proposition~\ref{prop:twosand} & \lean{le\_of\_ceiling\_le\_floor} & OLG/\allowbreak Certainty\allowbreak Equivalent\\
Proposition~\ref{prop:cechain} & \lean{mul\_le\_mul\_of\_euler\_pair} & OLG/\allowbreak Certainty\allowbreak Equivalent\\
The damped multiplier & \lean{damped\_le\_one} & OLG/\allowbreak Certainty\allowbreak Equivalent\\
The gradient sum, and its warning & \lean{sum\_mul\_rpow\_pred\_le} & Analysis/\allowbreak Power\allowbreak Mean\\
\bottomrule
\end{longtable}
\endgroup


\end{document}
